\documentclass[12pt,a4paper,onecolumn]{article}
\title{Digest of The New York Times}
\author{The New York Times}
% \definecolor{dark-green}{rgb}{0,0.25,0}
% \usepackage[utf8]{inputenc}

\usepackage{fontspec,xunicode,xltxtra}
\setmainfont[Ligatures=Required,Mapping=tex-text]{Georgia}
% \setmainfont[Ligatures=Required,Mapping=tex-text]{Palatino Linotype}
\setsansfont[Ligatures=Required,Mapping=tex-text]{Myriad Pro}
\setmonofont[Ligatures=Required,Mapping=tex-text]{Courier New}
\usepackage{xeCJK}
\setCJKmainfont{宋体}
\setCJKmonofont{黑体}

% \usepackage{indentfirst}


\usepackage[dvips,dvipsnames,svgnames]{xcolor}
\definecolor{light-gray}{gray}{0.95}

%生成PDF的链接
\usepackage[dvipdfmx]{hyperref}
\hypersetup{
    bookmarks=true,         % show bookmarks bar?
    bookmarksopen=true,
    pdfpagemode=UseNone,    % options: UseNode, UseThumbs, UseOutlines, FullScreen
    pdfstartview=FitH,
    pdfborder=1,
    pdfhighlight=/P,
    pdfauthor={wuxch},
    unicode=true,           % non-Latin characters in Acrobat’s bookmarks
    colorlinks,             % false: boxed links; true: colored links
    linkcolor=blue,         % color of internal links
    citecolor=blue,        % color of links to bibliography
    filecolor=magenta,      % color of file links
    urlcolor=cyan           % color of external links
}
% 使用目录
\usepackage{makeidx}
% 使用图形
\usepackage{graphicx}


\usepackage{titletoc}
\dottedcontents{section}[1.5em]{}{2.3em}{1pc}

% % 页眉页脚的设定
% \usepackage{fancyhdr}
% \pagestyle{fancy}
% \lhead{\includegraphics[height=0.2in]{zte_logo}}

% \usepackage{footmisc}
\usepackage{pfnote}
\usepackage{fnpos}

% % 嵌入的代码显示
% \usepackage{listings}
% \lstset{language=C++, breaklines, extendedchars=false}
% \lstset{keywordstyle=\color{blue},
%         commentstyle=\color{SeaGreen},
%         stringstyle=\ttfamily,
%         showstringspaces=false,
%         backgroundcolor=\color{light-gray}}
% \lstset{basicstyle=\small}

\usepackage{booktabs}

\usepackage{tipa}
\usepackage{phonetic}
\usepackage[gen]{eurosym}

% \usepackage{CJK}
\usepackage{makeidx}
% \usepackage{glossary}

\usepackage{fancyhdr}

\usepackage{lastpage}
\cfoot{\textsf{\thepage\ of \pageref*{LastPage}}}

\pagestyle{fancy}
\fancyhead{}
\renewcommand{\headrulewidth}{0pt}
\pagenumbering{arabic}

\usepackage{graphicx}
\usepackage{picinpar}

\setlength{\parindent}{0in}
\setlength{\parskip}{2.5ex plus4mm minus3mm}
% \setlength{\parskip}{.3\baselineskip}

\setcounter{section}{0}

% bibtex
\usepackage[super]{natbib}
% \usepackage[globalcitecopy,labelstoglobalaux]{bibunits}
% \makeatletter
% \renewcommand{\@bibunitname}{\jobname.\the\@bibunitauxcnt}
% \makeatother

% \renewcommand{\refname}{缩略语}
% \renewcommand{\refname}{}

\makeatletter                   % Reference list option change
\renewcommand\@biblabel[1]{#1:  } % from [1] to 1: 
\makeatother                    %

\newcommand{\citenumfont}[1]{\textsf{#1}}

% \usepackage{tipa}
% \usepackage{phonetic}

\newcommand{\phonetic}[1]{{\fontspec[Colour=CC00CC]{Microsoft Sans Serif} #1}}


%%% Local Variables:
%%% TeX-master: "medical_trivia"
%%% End:

\makeindex
\begin{document}
\date{}
% \maketitle
\tableofcontents
\pagebreak
\section{Is Facebook Growing Up Too Fast?}

WHEN Facebook signed up its 100 millionth member last August, its employees spread out in two parks
in Palo Alto, Calif., for a huge barbecue. Sometime this week, this five-year-old start-up, born in
a dorm room at Harvard, expects to register its 200 millionth user.

That staggering growth rate -- doubling in size in just eight months -- suggests Facebook is rapidly
becoming the Web's dominant social ecosystem and an essential personal and business networking tool
in much of the wired world.

Yet Facebook executives say they aren't planning to observe their latest milestone in any
significant way. It is, perhaps, a poor time to celebrate. The company that has given users new ways
to connect and speak truth to power now often finds itself as the target of that formidable
grass-roots firepower -- most recently over controversial changes it made to users' home pages.

As Facebook expands, it's also struggling to match the momentum of hot new start-ups like Twitter,
the micro-blogging service, while managing the expectations of young, tech-savvy early adopters,
attracting mainstream moms and dads, and justifying its hype-carbonated valuation.

By any measure, Facebook growth is a great accomplishment. The crew of Mark Zuckerberg, the
company's 24-year-old co-founder and chief executive, is signing up nearly a million new members a
day, and now more than 70 percent of the service's members live overseas, in countries like Italy,
the Czech Republic and Indonesia. Facebook's ranks in those countries swelled last year after the
company offered its site in their languages.

All of this mojo puts Facebook on a par with other groundbreaking -- and wildly popular -- Internet
services like free e-mail, Google, the online calling network Skype and e-commerce sites like eBay.
But Facebook promises to change how we communicate even more fundamentally, in part by digitally
mapping and linking peripatetic people across space and time, allowing them to publicly share myriad
and often very personal elements of their lives.

Unlike search engines, which ably track prominent Internet presences, Facebook reconnects regular
folks with old friends and strengthens their bonds with new pals -- even if the glue is nothing more
than embarrassing old pictures or memories of their second-grade teacher.

Facebook can also help rebuild families. Karen Haber, a mother of two living outside Tel Aviv, logs
onto Facebook each night after she puts the children to bed. She searches for her family's various
surnames, looking for relatives from the once-vast Bachenheimer clan of northern Germany, which
fractured during the Holocaust and then dispersed around the globe.

Among the three dozen or so connections she has made on Facebook over the last year are a fifth
cousin who is a clinical social worker in Woodstock, N.Y.; a fourth cousin running an eyeglasses
store in Zurich; and another fifth cousin, living in Hong Kong selling diamonds. Now she shares
memories, photographs and updates with them.

``I was never into genealogy and now suddenly I have this tool that helps me find the descendants of
people that my grandparents knew, people who share the same truth I do,'' Ms.~Haber says. ``I'm
using Facebook and trying to unite this family.''

Facebook has also become a vehicle for broad-based activism -- like the people who organized on the
site last year and mobilized 12 million people to march in protests around the globe against
practices of the FARC rebels in Colombia.

Discussing Facebook's connective tissue, Mr.~Zuckerberg recalls the story of Claus Drachmann, a
schoolteacher in northern Denmark who became a Facebook friend of Anders Fogh Rasmussen, Denmark's
prime minister. Mr.~Drachmann subsequently invited Mr.~Rasmussen to speak to his class of
special-needs children; the prime minister obliged last fall.

Mr.~Zuckerberg says the story illustrates Facebook's power to cut through arbitrary social barriers.
``This represents a generational shift in technology,'' he says. ``To me, what is interesting was
that it was possible for a regular person to reach the prime minister and that that interaction
happened.''

As Facebook has matured, so has Mr.~Zuckerberg. He has recently traded his disheveled, unassuming
image for an ever-present tie and making visits to media outfits like ``The Oprah Winfrey Show.''
And he says Facebook's most important metrics are not its membership but the percentage of the wired
world that uses the site and the amount of information -- photographs, news articles and status
updates -- zipping across its servers.

Facebook's mission, he says, is to be used by everyone in the world to share information seamlessly.
``Two hundred million in a world of six billion is tiny,'' he says. ``It's a cool milestone. It's
great that we reached that, especially in such a short amount of time. But there is so much more to
do.''

AS Facebook stampedes along, it still has to get out of its own way to soothe the injured feelings
of users like Liz Rabban.

Ms.~Rabban, 40, a real estate agent and the mother of two from Livingston, N.J., joined the site in
November 2007, quickly amassing 250 friends and spending hours on the site each day.

But these days, she spends less time on the site and posts caustic comments about Facebook's new
design, which turns a majority of every user's home page into a long ``stream'' of recent, often
trivial, Twitter-like updates from friends.

``The changes just feel very juvenile,'' Ms.~Rabban says. ``It's just not addressing the needs of my
generation and my peers. In my circle, everyone is pretty devastated about it.''

Ms.~Rabban is not alone. More than two and a half million dissenters have joined a group on
Facebook's own site called ``Millions Against Facebook's New Layout and Terms of Service.'' Others
are lambasting the changes in their own status updates, which are now, ironically, distributed much
more visibly to all of their Facebook friends.

The changes, Facebook executives say, are intended to make the act of sharing -- not just
information about themselves but what people are doing now -- easier, faster and more urgent. Chris
Cox, 26, Facebook's director of products and a confidant of Mr.~Zuckerberg, envisions users
announcing where they are going to lunch as they leave their computers so friends can see the
updates and join them.

``That is the kind of thing that is not meaningful when it is announced 40 minutes later,'' he says.

The simmering conflict over the design change speaks to the challenges of pleasing 200 million
users, many of whom feel pride of ownership because they helped to build the site with free labor
and very personal contributions.

``They have a strange problem,'' says S.~Shyam Sundar, co-director of the Media Effects Research
Laboratory at Pennsylvania State University, of Facebook's quandary. ``This is a technology that has
inherently generated community, and it has gotten to the point where members of that community feel
not only vested but empowered to challenge the company.''

Those tensions boiled up previously, when Facebook announced the intrusive Beacon advertising system
in 2007, and again when Facebook introduced new service terms earlier this year, which appeared to
give the company broad commercial control over the content people uploaded to the site.

Facebook responded to protests over the second move by promising users a vote in how the site would
be governed.

But while Facebook is willing to give users a voice, it doesn't necessarily want to listen.

Users are widely opposed to terms that grant Facebook the right to license, copy and disseminate
members' content worldwide. But Facebook says it has to ignore those objections to protect itself
against lawsuits from users who might blame the company if they later regret having shared some
piece of information with their friends. (Other Web sites have similar stipulations.)

While Facebook addressed the feedback on its unpopular design changes last week -- partly by saying
it would give users more control over the stream of updates that appear on their pages -- it also
said members' pages would soon become even busier and more dynamic, updating automatically instead
of requiring users to refresh their browsers to see new posts.

That's a change that may irk users like Ms.~Rabban, who don't like how busy their pages have become.
Facebook executives counter that it will help users share more information, and that they will
eventually come to appreciate it, just as they have with previous changes that were initially
jarring.

``It's not a democracy,'' Mr.~Cox says of his company's relationship with users. ``We are here to
build an Internet medium for communicating and we think we have enough perspective to do that and be
caretakers of that vision.''

PEOPLE, of course, sometimes like to keep secrets and maintain separate social realms -- or at least
a modicum of their privacy. But Facebook at almost 200 million members is a force that reinvents and
tears at such boundaries. Teachers are yoked together with students, parents with their children,
employers with their employees.

Uniting disparate groups on a single Internet service runs counter to 50 years of research by
sociologists into what is known as ``homophily'' -- the tendency of individuals to associate only
with like-minded people of similar age and ethnicity.

Facebook's huge growth is creating inevitable collisions as the whole notion of ``friend'' takes on
a highly elastic meaning. When the Philadelphia Eagles allowed the star safety Brian Dawkins to
leave for the Denver Broncos earlier this month, Dan Leone, a gate chief at Lincoln Financial Field,
the Eagles' stadium, expressed his disappointment by referring to the situation with an obscenity on
his Facebook status update.

Mr.~Leone's boss, who was his Facebook friend, forwarded the update to an Eagles guest services
manager, who fired him. The team has since refused to reconsider the matter, despite Mr.~Leone's
deep remorse and his star turn on countless radio talk shows across the country to discuss the
situation.

``If you know your boss is online, or anyone close to your boss is online, don't be making comments
that can be detrimental to your employment,'' Mr.~Leone advises.

Facebook is trying to teach members to use privacy settings to manage their network so they can
speak discreetly only to certain friends, like co-workers or family members, as opposed to other
``friends'' like bosses or professional colleagues. But most Facebook users haven't taken advantage
of the privacy settings; the company estimates that only 20 percent of its members use them.

Other problems are trickier, especially among true friends and family members. How, for example, can
Facebook remain a place for teenagers to share what they did on Saturday night when it is also the
place where their parents are swapping investment tips with old friends?

In the six weeks since Rich Hall, a 52-year-old theater manager in Mount Carroll, Ill., joined
Facebook, he has reconnected with more than 400 friends and acquaintances, including former high
school friends, his auto mechanic and former buddies from his days as a stock car driver.

In the course of his new half-hour-a-day Facebook habit, Mr.~Hall also ``friended'' the 60 high
school students he is directing in a school play, so he could coordinate rehearsal times. That led
some of them to deny his request because, as he says they told him, their parents ``found it
creepy.'' Along the way, Mr.~Hall also found photographs of his 19-year-old son on the site,
drinking beer at a Friday night bonfire.

``He denied it and said he wasn't there,'' Mr.~Hall says. ``I said, 'Let's go to this page together
and look at these photos.' Of course he did it. There are no secrets anymore.''

Dwindling secrets, and prying eyes, are at the heart of the Facebook conundrum. While offering an
efficient and far-reaching way for people to bond, the site has also eroded sometimes natural
barriers.

``People usually spend a lot of time trying to be separate -- parents and children are a good
example,'' says Danah Boyd, a social scientist who has studied social networks and now works in the
research department of Microsoft, which has invested in Facebook. ``You are already seeing young
people sitting there thinking, 'Why am I hanging out with my mother who is reminiscing with her high
school mates?' You are seeing some reticence with young people that wasn't there two years ago.''

For their part, Facebook executives say they are less interested in being cool than in being a
useful place where anyone can go to share elements of their lives.

``The people who started the company weren't cool. I'm not cool,'' Mr.~Cox says. ``If you look at
the people who work here, it's much more nerdy and curious than cool.

``Cool only lasts for so long, but being useful is something that applies to everyone.''

MR. ZUCKERBERG hopes that being ubiquitous and useful translates to the bottom line.

Though Facebook is privately held and doesn't publicly disclose its earnings, various press and
analysts' estimates of its 2008 revenues span from \$250 million to \$400 million. That range may
not be enough to cover the company's escalating expenses, and it hardly justifies some of the
atmospheric valuations that have been placed on the start-up, including the \$15 billion that
Microsoft assigned to the company when it invested in it in 2007.

Facebook's financial challenges aren't unique. Popular free e-mail services like Hotmail from
Microsoft and Gmail from Google have little in the way of profits to show for their vast audiences,
aside from a few text ads that people rarely click on. Instant messaging networks like Microsoft
Messenger and AIM from American Online are similarly popular but have never been hyperprofitable,
for the simple reason that people do not want intrusive ads inserted into personal conversations.

Facebook's approach is to invite advertisers to join in the conversation. New ``engagement'' ads ask
users to become fans of products and companies -- sometimes with the promise of discounts. If a
person gives in, that commercial allegiance is then broadcast to all of the person's friends on the
site.

A new kind of engagement ad, now being tested, will invite people to vote -- ``what's your favorite
color M\&M?'' for example -- and brands will pay every time a Facebook member participates.

``We are trying to provide the antidote for the consumer rebellion against interruptive
advertising,'' says Sheryl Sandberg, Facebook's chief operating officer and Mr.~Zuckerberg's
business consigliere.

Ms.~Sandberg, who ran Google's highly successful advertising initiatives before leaving the search
giant to join Facebook, said her company's revenue was growing despite a brutal downturn that is
hurting other kinds of online advertising. She also puts one rumor to rest, saying the company is
not considering charging members for any aspect of its service.

``We're pretty pleased with the overall trajectory,'' she says. ``Our conversations with big
advertisers have broadened in scope and we also have more people asking about how they can work with
us.''

Facebook recently introduced advertising tools to let companies focus on users based on the language
they use on the site and their geographic location. So, for example, an advertiser can now tailor a
message to the Latino community in Los Angeles or French speakers in Montreal.

Despite the gloom permeating much of the advertising world, and the formidable challenges facing the
site, some advertisers say they glimpse the future in Facebook's brand of interactive advertising.

``Our clients all want to see if they can make this work,'' says Al Cadena, the interactive account
director at Threshold Interactive in Los Angeles, which represents companies like Nestl\'e, Honda
and Sony. ``Advertising used to be a one-way communication from advertiser to consumer, but now
people want to have a dialogue. And Facebook is becoming the default way to do that, not only in the
States but really for the whole world.''

Internet evangelists say that when a technology diffuses into society, as Facebook appears to be
doing, it has achieved ``critical mass.'' The sheer presence of all their friends, family and
colleagues on Facebook creates potent ties between users and the site -- ties that are hard to break
even when people want to break them.

Many who have tried to free themselves of their daily Facebook habit and leave the site, like Kerry
Docherty, a student at Pepperdine University's law school, speak of a powerful gravitational pull
and an undercurrent of peer pressure that eventually brings them back.

``People gave me a hard time for leaving Facebook,'' says Ms.~Docherty, who quit at the end of 2007
but then rejoined six months later. ``Everyone has a love-hate relationship with it. They wanted me
to be wasting my time on it just like they were wasting their time on it.''

\section{G.M.~Chief Quits After Pressure From Administration}

DETROIT -- The chairman and chief executive of General Motors, Rick Wagoner, resigned Sunday as part
of a broad agreement with the Obama administration to funnel more government aid to the ailing auto
giant, according to people close to the decision.

Mr.~Wagoner, who has served as G.M.'s top executive since 2000, agreed to step down after it was
requested by the president's auto task force, these people said.

G.M. had no immediate comment on the sudden development, which came on the eve of Mr.~Obama's
announcement on Monday that is to detail his rescue plans for G.M., Chrysler and the larger American
auto industry.

But people in the company said G.M. would issue a statement on Mr.~Wagoner after the president
unveiled his plan in Washington.

``The bigger surprise is not that he resigned. That was going to happen sooner or later,'' said
Michael Useem, a professor of management at the Wharton School. ``But the moment seems
inexplicable.''

The president's task force is expected to recommend more short-term aid for G.M. and Chrysler, but
with tight strings on the money and a deadline on getting concessions from union workers and
creditors.

A person with direct involvement in the auto bailout discussions said the new deadline would be
April 30.

``Thirty days from now, there will either be a bankruptcy or the naming of a chief restructuring
officer who will have government authority to `knock heads together,' '' this person said of G.M. In
addition, the government must come up with a backup guarantee on loan for G.M. to operate in
bankruptcy because the banks will not do it.

G.M. and Chrysler have almost exhausted the \$17.4 billion in federal aid the two companies have
received since December. G.M. has asked for up to an additional \$16.6 billion, and Chrysler has
requested an additional \$5 billion.

According to people close to the talks, the task force will treat G.M. and Chrysler differently with
respect to their overhlaul plans and aid requests.

Chrysler has submitted a stand-alone revival plan, but has also proposed entering a global alliance
with the Italian automaker Fiat.

People with knowledge of the task force's deliberations said that members so far are looking
favorably on the deal, which would give Fiat a 35 percent stake in Chrysler in exchange for
providing small cars and engines to the American company.

Like Chrysler, G.M. submitted an overhaul plan in February, which called for cutting 47,000 jobs
worldwide and drastically shrinking the company's models, brands and dealers.

But people close to the auto task force said some members had considered the plan inadequate to
transform G.M. into a profitable enterprise.

Mr.~Obama, in comments in a televised interview on Sunday, said neither G.M. nor Chrysler had yet
met the conditions of their existing loans.

``That's going to mean a set of sacrifices from all parties involved -- management, labor,
shareholders, creditors, suppliers, dealers. Everybody's going to have to come to the table and say
it's important for us to take serious restructuring steps now in order to preserve a brighter future
down the road,'' Mr.~Obama said in a taped interview on the CBS news program ``Face the Nation.''

While Mr.~Obama did not specify a need to replace Mr.~Wagoner at G.M., the president has repeatedly
cited mistakes made by management as contributing to the industry's troubles.

Mr.~Wagoner was not available for comment, according to people at G.M. As recently as March 18 he
said in an interview that he did not consider his job at stake in his discussions with the
president's auto task force. ``They so far haven't commented on that,'' he said.

Administration officials stopped short of saying that Mr.~Wagoner had been forced out, only that he
was asked to leave and agreed, according to people with knowledge of the decision.

Frederick A.~Henderson, G.M.'s president, is a leading candidate to replace Mr.~Wagoner as CEO, at
least temporarily. Members of the auto panel spoke with Mr.~Henderson recently, and came away with a
favorable impression of him, people with knowledge of the panel's discussions said Sunday. But the
task force planned to recruit an outside reorganization specialist to run the company over the long
term, one person said.

Mr.~Wagoner's departure at G.M. marks an end to a corporate hierarchy that spanned generations. The
only previous G.M. chairman to leave under duress was Robert C.~Stempel, who was force to resign in
1992.

Mr.~Wagoner, a graduate of Duke University and the Harvard Business School, vaulted into Detroit's
consciousness in 1992, when he was named G.M.'s chief financial officer at the age of 38.

He was a prot\'eg\'e of G.M.'s former chief executive, John F.~Smith Jr., who ran the company after
Robert C.~Stempel resigned in October 1992. Mr.~Wagoner stepped in when G.M.'s purchasing chief,
Ignacio Lopez de Arriortua, left for Volkswagen in 1993.

A year later, Mr.~Wagoner was named president of G.M.'s North American operations, and was elevated
to president of the company and its chief operating officer in 1998. He succeeded Mr.~Smith as chief
executive in 2000, and became G.M.'s chairman in 2003.

Mr.~Wagoner followed Mr.~Smith in expanding G.M.'s operations outside the United States. In 2007,
G.M. sold more vehicles outside North American than it did in its core market, in part because of
Mr.~Wagoner's aggressive pursuit of sales in China, Latin America and Eastern Europe.

But G.M.'s share of its most important market, the United States, declined steadily under
Mr.~Wagoner. In 1994, when he took charge of North America, G.M. held 33.2 percent of the American
market. Last month, G.M. held only 18.8 percent of American auto sales, according to statistics from
Motorintelligence.com, which specializes in industry data. Auto sales in February were the worst for
the industry since 1981.

Shortly after taking the chief executive's job, Mr.~Wagoner predicted that G.M. could earn as much
as \$10 a share by the middle of the decade, because of efforts the company expected to take to cut
its costs, transform its operations and build its business.

Last year, Mr.~Wagoner led the celebration for G.M.'s 100th birthday, promising to steer the
automaker into its next century with new technology and a renewed vigor. But G.M. collapsed last
fall when new-vehicle sales in the United States plummeted to their lowest level in 25 years. G.M.
lost more than \$30 billion in 2008, and has been subsisting on government loans since the beginning
of the year.

As G.M.'s biggest defender over the years, Mr.~Wagoner has long been the target of critics including
shareholders.

``He's a victim of problems that his predecessors did not solve, but he's also responsible for where
G.M. is today,'' Mr.~Useem said.


\section{U.S. Moves to Overhaul Ailing Carmakers}

WASHINGTON -- The White House on Sunday pushed out the chairman of General Motors and instructed
Chrysler to form a partnership with the Italian automaker Fiat within 30 days as conditions for
receiving another much-needed round of government aid.

The decision to ask G.M.'s chairman and chief executive, Rick Wagoner, to resign caught Detroit and
Washington by surprise, and it underscored the Obama administration's determination to keep a tight
rein on the companies it is bailing out -- a level of government involvement in business perhaps not
seen since the Great Depression.

President Obama is scheduled to announce details of the auto package at the White House on Monday,
but two senior officials, offering a preview on condition of anonymity, made clear that some form of
bankruptcy -- a quick, court-supervised restructuring, as they described it -- could still be an
option for one or both companies.

Mr.~Obama's auto industry task force, in a report released Sunday night assessing the viability of
both companies and detailing the administration's new plans for them, concluded that Chrysler could
not survive as a stand-alone company.

The report said the company would get no more help from the government unless it can finalize a
proposed alliance with the Italian automaker Fiat by April 30. It must also reduce its debt and
health-care obligations.

If a deal is reached between Chrysler and Fiat, the administration says it would consider another
loan of \$6 billion to Chrysler.

G.M., on the other hand, has made considerable progress in developing new energy-efficient cars and
could survive if it can cut costs sharply, the task force reported. The administration is giving
G.M. 60 days to present a cost-cutting plan and will provide taxpayer assistance to keep it afloat
during that time.

Along with Mr.~Wagoner's ouster, the task force said most of the company's board of directors would
be replaced over the next few months.

Mr.~Wagoner's resignation is the latest example of the government taking a hands-on role in making
major decisions at companies it is bailing out. The government has already pushed banks to make
management changes and sharply reduce or eliminate their dividends, and it also is directing many of
the decisions at the troubled insurance giant American International Group, which is nearly 80
percent owned by the government after its rescue.

In deciding to urge Mr.~Wagoner to step down, the Obama administration seemed mindful of the
public's growing outrage over bailouts of private companies, as well as the bonuses paid to
employees of A.I.G.

Mr.~Obama is well aware that he cannot afford to give the appearance of using tax dollars to reward
executives who have done a poor job, and he began signaling as early as last week that he would take
a tough stance with the automakers.

In a question and answer session at the White House on Thursday, the president said there had been
``a lot of mismanagement of the auto industry over the past several years,'' and declared that more
government help would be contingent on the companies' ``willingness to make some pretty drastic
changes.''

The plan Mr.~Obama is to announce on Monday will also include government backing of warranties for
G.M. and Chrysler cars and trucks, to give consumers enough confidence to buy them, even if one or
both are forced into bankruptcy.

In Detroit, G.M. had no immediate comment on the sudden resignation of Mr.~Wagoner. Company
officials said G.M. would issue a statement about Mr.~Wagoner after Mr.~Obama revealed his plan in
Washington.

Mr.~Wagoner has presided over a steep drop in G.M.'s domestic market share, which has led to tens of
billions of dollars in losses. His critics have said that management's failure to move aggressively
to address the company's problems contributed to its dire financial situation.

``The bigger surprise is not that he resigned. That was going to happen sooner or later,'' said
Michael Useem, professor of management at the Wharton School of Business at the University of
Pennsylvania. ``But the moment seems inexplicable.''

G.M. and Chrysler have almost exhausted the combined \$17.4 billion in federal aid they have
received since December. G.M. has asked for up to \$16.6 billion more, and Chrysler has requested
another \$5 billion.

Bondholders are under pressure to convert two-thirds of the \$27 billion owed them into G.M. stock,
while the United Auto Workers union is being asked to substitute stock for 50 percent of their
health care benefits for retirees. Both groups have resisted those changes.

Administration officials say they have enough money to offer the assistance they envision under
plans already approved by Congress. Even so, Mr.~Obama may face skepticism on Capitol Hill and from
the public.

As part of the companies' original agreement for the loans, both were required to submit
restructuring plans. Mr.~Wagoner's removal underscores how much more G.M. needs to cut than was
proposed in the plan the company submitted.

Administration officials stressed that the company needed a fresh approach and leadership changes;
they said Steven Rattner, the former investment banker who co-chairs the auto task force, delivered
the news to Mr.~Wagoner.

People at G.M. said Mr.~Wagoner was not available for comment Sunday night. As recently as March 18
he said in an interview that his discussions with the task force did not give him the impression
that his job was at stake. ``They so far haven't commented on that,'' he said then.

Frederick A.~Henderson, G.M.'s president, will replace Mr.~Wagoner on an interim basis as chief
executive; Kent Kresa, a board member, will assume the chairmanship. Members of the auto panel spoke
with Mr.~Henderson recently and came away with a favorable impression of him, people familiar with
the panel's discussions said.

Like Mr.~Wagoner, Mr.~Henderson is a graduate of the Harvard Business School and a lifer at G.M. He
started in the finance division in 1984 and later spent nine years in executive positions in South
America, Asia and Europe. The Detroit-born son of a G.M. sales manager, Mr.~Henderson, 50, became
chief financial officer in 2006 and was named president and chief operating officer a year ago.

Mr.~Wagoner's departure at G.M. marks an end to a corporate hierarchy that spanned generations. The
last G.M. chairman to leave under duress was Robert C.~Stempel, who was forced out in 1992 by
outside directors who blamed him for losses.

Mr.~Wagoner, 56, came to G.M. in 1977 and rose to become chief financial officer in 1992 when he was
38. He oversaw the company's North American business for years before being named chairman in 2000.

G.M.'s share of its most important market, the United States, declined steadily under Mr.~Wagoner.
In 1994, when he took charge of North America, G.M. held 33.2 percent of the American car market.
Last month, G.M.'s share was only 18.8 percent, according to statistics from Motorintelligence.com,
which specializes in industry data. Auto sales in February were the worst for the industry since
1981.

G.M. collapsed last fall when new-vehicle sales in the United States plummeted to their lowest level
in 25 years. G.M. lost more than \$30 billion in 2008, and has been subsisting on government loans
since the beginning of the year.

The administration briefed lawmakers on the plan sunday night. Afterward, Representative Thaddeus
G.~McCotter, Republican of Michigan, whose district is just outside Detroit, expressed frustration
over the ousting of Mr.~Wagoner and with administration officials for not being clearer about the
potential job losses that lie ahead.

``Why would you ask Rick Wagoner to resign when you are giving G.M. 60 days to meet a new target,
but you aren't saying what the new goal is yet,'' Mr.~McCotter said in an interview.

\section{In Rescue of Captain, Navy Kills 3 Pirates}

Navy Seal snipers rescued an American cargo ship captain unharmed and killed three Somali pirates in
a daring operation in the Indian Ocean on Sunday, ending a five-day standoff between United States
naval forces and a small band of brigands in a covered orange lifeboat off the Horn of Africa.

Acting with President Obama's authorization and in the belief that the hostage, Capt.~Richard
Phillips, was in imminent danger of being killed by captors armed with pistols and AK-47s, snipers
on the fantail of the destroyer Bainbridge, which was towing the lifeboat on a 100-foot line, opened
fire and picked off the three captors.

Two of the captors had poked their heads out of a rear hatch of the lifeboat, exposing themselves to
clear shots, and the third could be seen through a window in the bow, pointing an automatic rifle at
the captain, who was tied up inside the 18-foot lifeboat, senior Navy officials said.

It took only three remarkable shots -- one each by snipers firing from a distance at dusk, using
night-vision scopes, the officials said. Within minutes, rescuers slid down ropes from the
Bainbridge, climbed aboard the lifeboat and found the three pirates dead. They then untied Captain
Phillips, ending the contretemps at sea that had riveted much of the world's attention. A fourth
pirate had surrendered earlier.

Shortly after his rescue, Captain Phillips was taken aboard the Bainbridge, underwent a medical exam
and was found to be in relatively good condition for a 53-year-old seafarer who had been held since
Wednesday by pirates who had demanded \$2 million for his life. He called home and was flown to the
Boxer, an amphibious assault ship also off the Somali coast. Arrangements were being made Sunday
night for his return home to Vermont.

``I share the country's admiration for the bravery of Captain Phillips and his selfless concern for
his crew,'' Mr.~Obama said in a statement. ``His courage is a model for all Americans.''

Jubilation over the dramatic rescue reached from the White House to Underhill, Vt., Captain
Phillips's hometown, and from personnel aboard the Bainbridge to the cheering, fist-pumping
19-member crew of the captain's cargo ship, the Maersk Alabama, docked in Mombasa, Kenya.

Captain Phillips, who was said to be resting comfortably, spoke to officials of the Maersk Line, who
quoted him as saying: ``The real heroes are the Navy, the Seals, those who have brought me home.''
He also spoke to his wife, Andrea, and two college-aged children in Underhill, where dozens of
yellow ribbons fluttered on the white picket fence of his home and two small American flags jutted
up from the lawn.

``This is truly a very happy Easter for the Phillips family,'' said Alison McColl, a Maersk
representative assigned to speak for the family. ``They are all just so happy and relieved,'' she
said. ``I think you can all imagine their joy and what a happy moment it was for them.''

On the family's behalf, Ms.~McColl thanked the nation and the people of Vermont for their prayers
and support. ``Obviously, this has been a long journey for the family,'' she said. John Reinhart,
president and chief executive of Maersk Line Ltd., praised the Navy and federal officials for their
performance. ``Everyone's worked around the clock,'' he said. ``It's magnificent to see the
outcome.''

While the outcome was a triumph for America, officials in many countries plagued by pirates said it
was not likely to discourage them. Pirates are holding a dozen ships with more than 200 crew
members, according to the Malaysia-based International Maritime Bureau.

In Somalia itself, other pirates reacted angrily to the news that Captain Phillips had been rescued,
and some said they would avenge the deaths of their colleagues by killing Americans in sea
hijackings to come.

``Every country will be treated the way it treats us,'' Abdullahi Lami, one of the pirates holding a
Greek ship anchored in the pirate den of Gaan, a central Somali town, was quoted by The Associated
Press as saying in a telephone interview. ``In the future, America will be the one mourning and
crying.''

Aboard the Maersk Alabama, a 17,000-ton cargo ship, Captain Phillips's crew erupted in cheers, waved
American flags and fired off flares. When four pirates attacked the ship on Wednesday, the crew
escaped harm after the captain offered himself as a hostage. He told his crewmen to lock themselves
in cabins, and allowed himself to be taken at gunpoint into the lifeboat in which the pirates fled.

Over the ensuing days, according to official accounts of the episode, the pirates made repeated
threats to kill the captain as their motorized lifeboat moved about 30 miles off the Somali coast.
It was closely watched by United States warships and helicopters in an increasingly tense standoff.

Talks to free the captain began Thursday, with the commander of the Bainbridge communicating with
the pirates under instructions from F.B.I. hostage negotiators flown to the scene. The pirates
threatened to kill Captain Phillips if attacked, and the result was tragicomic: the world's most
powerful navy vs. a lifeboat.

Vice Adm. William E.~Gortney, commander of the United States naval forces in the region, said in a
briefing in Bahrain that despite ransom demands from the pirates the United States had not discussed
any ransom and had talked to the pirates only about the release of Captain Phillips and the pirates'
surrender.

The Defense Department twice sought Mr.~Obama's permission to use force to rescue Captain Phillips,
most recently on Friday night, senior defense officials said. On Saturday morning, the president
agreed, they said, if it appeared that the captain's life was in imminent danger.

By Friday, with several warships within easy reach of the lifeboat, the negotiations had gone
nowhere. Captain Phillips jumped into the sea, but was quickly recaptured. On Saturday, the pirates
fired several shots at a small boat that had approached from the Bainbridge.

By the weekend, however, the pirates had begun to run out of food, water and fuel. That apparently
provided the opening officials were hoping for. In briefings, senior officers who spoke anonymously
because they had not been authorized to disclose information said that the pirates agreed to accept
food and water. A small craft was used to deliver them and it apparently made several trips between
the Bainbridge and the lifeboat.

On one trip, one of the four pirates -- whose hand had been gashed during the capture of Captain
Phillips -- asked for medical treatment and, in effect surrendering, was taken in the small boat to
the Bainbridge. Justice Department officials were studying options for his case, including criminal
charges in the United States or turning him over to Kenya, where dozens of pirates have faced
prosecution. Three pirates were left on board with Captain Phillips.

Meanwhile, members of the Navy Seals were flown in by fixed-wing aircraft. They parachuted into the
sea with inflatable boats and were picked up by the Bainbridge. On Sunday, the pirates, their fuel
gone, were drifting toward the Somali coast. They agreed to accept a tow from the Bainbridge, the
senior officials said. At first, the towline was 200 feet long, but as darkness gathered and seas
became rough, the towline was shortened to 100 feet, the officials said. It was unclear if this was
done with the pirates' knowledge.

At dusk, a single tracer bullet was seen fired from the lifeboat. The intent was unclear, but it
ratcheted up the tension and Seal snipers at the stern rail of the Bainbridge fixed night-vision
scopes to their high-powered rifles, getting ready for action.

What they saw was the head and shoulders of two of the pirates emerging from the rear hatch of the
lifeboat. Through the window of the front hatch they saw the third pirate, pointing his AK-47 at the
back of Captain Phillips, who was seen to be tied up.

That was it: the provocation that fulfilled the president's order to act only if the captain's life
was in imminent danger, and the opportunity of having clear shots at each captor. The order was
given. Senior defense officials, themselves marveling at the skill of the snipers, said each took a
target and fired one shot.

``This was an incredible team effort,'' Admiral Gortney said when it was over. ``And I am extremely
proud of the tireless efforts of all the men and women who made this rescue possible.''

\section{Utility in the Jumble of Tweets}

SAN FRANCISCO -- The first reaction many people have to Twitter is befuddlement. Why would they want
to read short messages about what someone ate for breakfast?

It's a reasonable question. Twitter unleashes the diarist in its 14 million users, who visited its
site 99 million times last month to read posts tapped out with cellphones and computers.

Individually, many of those 140-character ``tweets'' seem inane.

But taken collectively, the stream of messages can turn Twitter into a surprisingly useful tool for
solving problems and providing insights into the digital mood. By tapping into the world's
collective brain, researchers of all kinds have found that if they make the effort to dig through
the mundane comments, the live conversations offer an early glimpse into public sentiment -- and
even help them shape it.

Companies like Starbucks, Whole Foods and Dell can see what their customers are thinking as they use
a product, and the companies can adapt their marketing accordingly. Last week in Moldova, protesters
used Twitter as a rallying tool while outsiders peered at their tweets to help them understand what
was happening in that little-known country.

And over the weekend, Amazon.com learned how important it was to respond to the Twitter audience.
After one author noticed that Amazon had reclassified books with gay and lesbian themes as ``adult''
and removed them from the main search and sales rankings, a protest broke out on blogs and Twitter.
The company felt compelled to respond despite the Easter holiday, initially saying the problem was
due to a ``glitch in our system'' but later blaming a ``ham-fisted cataloging error'' that affected
more than 57,000 books dealing with health and sex.

Soon, machines could twitter as much as people. Corey Menscher, a graduate student at New York
University, developed the Kickbee, an elastic band with vibration sensors that his pregnant wife
wore to alert Twitter each time the baby kicked: ``I kicked Mommy at 08:52 PM on Fri, Jan 2!''
Mr.~Menscher is now considering selling the product.

Pairing sensors with Twitter leads some to think Twitter could be used to send home security alerts
or tell doctors when a patient's blood sugar or heart rate climbs too high. In the aggregate, such
real-time data streams could aid medical researchers.

Already doctors use Twitter to ask for help and share information about procedures. At Henry Ford
Hospital in Detroit, surgeons and residents twittered throughout a recent operation to remove a
brain tumor from a 47-year-old man who has seizures.

``A portion of the skull is being removed to allow access to the dura, the lining of the brain,'' an
early tweet said. Medical residents and curious laymen following online asked the doctors what music
they were listening to (Loreena McKennitt, a Celtic singer), whether the patient felt pain in the
brain (no, just pressure) and how big the tumor was (the size of a golf ball). As is convention on
Twitter, they tagged all their tweets with a keyword so anyone could search for the keyword and read
the stream of posts.

``Twitter lets people know what's going on about things they care about instantly, as it happens,''
said Evan Williams, Twitter's chief executive and co-founder. ``In the best cases, Twitter makes
people smarter and faster and more efficient.''

Mr.~Williams, along with the other founders, Biz Stone and Jack Dorsey, first envisioned Twitter as
an easy way to stay in touch with people you already know.

In 2006, when Twitter was just starting, the three men felt a small earthquake in San Francisco.
They each reached for their phones to twitter about it and discovered tweets from others in the
city. At that moment, it dawned on them that Twitter might be most useful for something else -- a
frontline news report, not just for friends, but for anyone reading.

Indeed, the news-gathering promise of Twitter was most evident during the terrorist attacks in
Mumbai last November and when a jetliner landed in the Hudson River in January. People were
twittering from the scenes before reporters arrived.

The attention the service received helped it nearly double the number of new users in the last
month, making Twitter the third-largest online social network, behind Facebook and MySpace,
according to Compete, a Web analytics company.

``Twitter reverses the notion of the group,'' said Paul Saffo, the Silicon Valley futurist.
``Instead of creating the group you want, you send it and the group self-assembles.''

Martin Stoll first discovered this on a visit to New York, where he was searching for a comedy show.
Minutes after posting a query on Twitter, five people he did not know had recommended shows. People
who had signed up to follow Mr.~Stoll's tweets had his question delivered to their Twitter page or
cellphone, and others reading the live Twitter stream could also see it.

Mr.~Stoll, the founder of GoSeeTell Network, an online travel company, realized Twitter could be an
on-the-go, living guidebook for tourists. He created the Portland Twisitor Center, where thousands
of people ask where to find the best brunch spot or coffee house and receive instant responses from
the center's officials and anyone else who wants to answer them.

Corporations often use Twitter for sales pitches. Intuit, the maker of QuickBooks and TurboTax,
monitors Twitter for people writing about Mint, a personal finance Web site that competes with its
Quicken Online. Intuit then writes to them and offers its service.

Even small businesses find Twitter useful. For example, Mary F.~Jenn, of True Massage and Wellness
in San Francisco, twitters when masseuses have same-day openings in their schedules and offers
discounts. The spa is often fully booked within several hours.

But Twitter's most productive use has been for businesses that want to peer into the minds of their
customers, reading their immediate reactions to a product. Dell noticed customers complaining on
Twitter that the apostrophe and return keys were too close together on the Dell Mini 9 laptop. So
Dell fixed the problem on the Dell Mini 10.

At Starbucks, customers used to complain by leaving notes in suggestion boxes. Now they can also
post their complaints or suggestions on Twitter, where Brad Nelson, who writes the company's Twitter
updates, tracks what people are saying about Starbucks online.

Last month, rumors surfaced that Starbucks would not send coffee to troops in Iraq in protest of the
war. Mr.~Nelson shot them down, twittering, ``This is not true. Get the facts here,'' with a link to
Starbucks's refutation of the rumor.

Some developers are creating tools to help companies keep an eye on the buzz. Akshay Java, a
scientist at Microsoft, is trying to figure out a way to identify which experts are most influential
on given topics by automatically analyzing the content of their tweets and who is in their Twitter
network. Companies like Microsoft could use that information to figure out which twitterers they
should contact to create buzz about a new product.

However, for Twitter to be truly useful as a research tool, more people will have to start using it.
If it collected a more representative slice of what the world is thinking, Twitter could enable
academics and scientists to track epidemics, for instance.

To make that easier, Twitter will soon add a search box to the home page so users can search for
terms like ``earthquake'' or ``flu'' and get any tweets about those topics in their Twitter feeds.

To continue growing, Twitter will also need to earn meaningful revenue, which the two-year-old
company has yet to do. Twitter hopes to charge companies like Starbucks for features that help them
communicate with and learn more about their customers, the founders said.

As the company taps into the \$35 million it recently raised from two venture capital firms in
Silicon Valley -- on top of the \$20 million it previously raised -- Mr.~Williams sees evidence that
his service has started to find a more utility-minded, mainstream audience.

He points to people using it to find gasoline in Atlanta during a gas shortage last fall. ``It was
so far from sharing what you had for breakfast -- yet it only works because it's the same place
where people talk about breakfast,'' he said.

\section{Labor Groups Reach an Accord on Immigration}

The nation's two major labor federations have agreed for the first time to join forces to support an
overhaul of the immigration system, leaders of both organizations said on Monday. The accord could
give President Obama significant support among unions as he revisits the stormy issue in the midst
of the recession.

John Sweeney, president of the A.F.L.-C.I.O., and Joe T.~Hansen, a leader of the rival Change to Win
federation, will present the outlines of their new position on Tuesday in Washington. In 2007, when
Congress last considered comprehensive immigration legislation, the two groups could not agree on a
common approach. That legislation failed.

The accord endorses legalizing the status of illegal immigrants already in the United States and
opposes any large new program for employers to bring in temporary immigrant workers, officials of
both federations said.

``The labor movement will work together to make sure that the White House as well as Congress
understand that we speak about immigration reform with one voice,'' Mr. Sweeney said in a statement
to The New York Times.

But while the compromise repaired one fissure in the coalition that has favored broad immigration
legislation, it appeared to open another. An official from the United States Chamber of Commerce
said Monday that the business community remained committed to a significant guest-worker program.

``If the unions think they're going to push a bill through without the support of the business
community, they're crazy,'' said Randel Johnson, the chamber's vice president of labor, immigration
and employee benefits. ``There's only going to be one shot at immigration reform. As part of the
trade-off for legalization, we need to expand the temporary worker program.''

The common labor position is also unlikely to convince many opponents that an immigration overhaul
would not harm American workers. When Obama administration officials said last week that the
president intended to push Congress this year to take up an immigration bill that would include a
path to legal status for the country's estimated 12 million illegal immigrants, critics criticized
the approach as amnesty for lawbreakers.

``In our current economic crisis, Americans cannot afford to lose more jobs to illegal workers,''
said Representative Steve King, an Iowa Republican who sits on the House Judiciary subcommittee on
immigration. ``American workers are depending on President Obama to protect their jobs from those in
America illegally.''

The two labor federations have agreed in the past to proposals that would give legal status to
illegal immigrants. But in 2007 the A.F.L.-C.I.O. parted ways with the service employees and several
other unions when it did not support legislation put forth by the Bush administration because it
contained provisions for an expanded guest-worker program.

In the new accord, the A.F.L.-C.I.O. and Change to Win have called for managing future immigration
of workers through a national commission. The commission would determine how many permanent and
temporary foreign workers should be admitted each year based on demand in American labor markets.
Union officials are confident that the result would reduce worker immigration during times of high
unemployment like the present.

Mr.~Hansen, who is president of the United Food and Commercial Workers Union, said in an interview
that the joint proposal was a ``building block to go forward to get immigration reform up on the
agenda in Congress'' sometime this year.

Thousands of immigrant farm workers and other low-wage laborers come to the United States through
seasonal guest-worker programs that are subject to numerical visa limits and have been criticized by
employers as rigid and inefficient. Many unions oppose the programs because the immigrants are tied
to one employer and cannot change jobs no matter how abusive the conditions, so union officials say
they undercut conditions for American workers. Highly skilled foreign technology engineers and
medical specialists also come on temporary visas.

Advocates for immigrants said a unified labor movement could substantially bolster their position as
they push for legislation to restructure the ailing immigration system.

``It shows how important the issue is to the representatives of American workers,'' said Frank
Sharry, executive director of America's Voice, an advocate group.

A.F.L.-C.I.O. officials said they agreed with Change to Win leaders that, with more than seven
million unauthorized immigrants already working across the nation, legalizing their status would be
the most effective way to protect labor standards for all workers.

``We have developed a joint strategy with the approach framed around workers' rights,'' said Ana
Avenda\~{n}o, associate general counsel of the A.F.L.-C.I.O.

Labor leaders said that they would talk with other groups in coming weeks to nail down details of a
common position, and that they would then would work in Congress and with the Obama administration
to try to ensure that their proposal was part of any bill offered for debate.

Also supporting the compromise is Eliseo Medina, an executive vice president of the Service
Employees International Union, a member of Change to Win with hundreds of thousands of members who
are immigrants. The Change to Win federation was formed in 2005 with seven unions that broke away
from the A.F.L.-C.I.O.

The plan for a labor commission to monitor and control levels of worker immigration was developed
with help from Ray Marshall, a labor secretary under President Jimmy Carter. Over the past year,
Mr.~Marshall, at the request of the A.F.L.-C.I.O., has been consulting between the two federations
and with a variety of Hispanic organizations and advocate groups for immigrants.

``All these groups understand that one of the main reasons they lost before was that they were not
together,'' Mr.~Marshall said.

According to a list of principles the labor leaders will present on Tuesday, they are proposing a
``depoliticized,'' independent commission that ``can assess labor market needs on an ongoing basis
and -- based on a methodology to be approved by Congress -- determine the number of foreign workers
to be admitted for employment purposes.''

Mr.~Johnson, the Chamber of Commerce official, said, ``A commission doesn't get us there.''

Tamar Jacoby, president of ImmigrationWorks USA, a group that organizes businesses to support
comprehensive immigration legislation, agreed that employers would have many questions about the
approach.

``The question is, Will the commission work?'' Ms.~Jacoby said. ``Will it be adequately attuned to
and triggered by the labor market? A system that may -- or may not -- supply the workers that
business will need in the future after the recession will be a cause of great concern to
employers.''

\section{Glamour Dims as Hecklers Hit the Auto Show}

Just a year ago, working as a product presenter at an auto show was a pretty straightforward job.
You stood next to a vehicle, you called it a marvel of engineering, style and comfort and then you
fielded softball questions like, ``What does this baby cost?''

But that was before the bailout. Now that the government has helped General Motors and Chrysler
stave off bankruptcy with billions of dollars in loans, these companies are finding somewhat hostile
crowds at their exhibits. Which leads to scenes like the one on Friday at the New York auto show,
where a blond woman in a tight black dress stood on a rotating platform and pitched the sporty Dodge
Circuit, one of five electric cars that Chrysler is developing.

Donald Han, an accountant from Queens, sounded unmoved. ``Why now?'' he asked the woman, rather
curtly, once she had finished her patter. ``How come you've got to nearly go bankrupt before you
come out with a car like this?''

Long a glamorous showcase for carmakers, auto shows have lately become a place for buyers and
gawkers to vent. Few of the attendees at the Javits Center, where the New York show runs until
Sunday, will ever encounter a top executive from G.M. or Chrysler. But all of them get within
heckling range of the presenters and for some, that is good enough.

It does not seem to matter that these women -- they are nearly all women, most of them young and
attractive -- work part time for marketing firms and talent agencies that have contracts to run the
exhibits. Many know little about the car companies they are working for beyond the scripts they have
memorized.

``I try to explain that we're not involved in corporate decisions, so complaining to us doesn't
really make a lot of sense,'' said Kerri Moss, standing on a large turntable next to a Jeep 4X4
Laredo, a Chrysler product. Recently laid off from her job as a teacher, she is trying to earn some
money on the car show circuit, which runs from September to May. ``And if that doesn't work, I tell
them we're doing the best we can.''

Often, that does not work either. One G.M. presenter said a woman told her the company was
responsible for the death of American soldiers in Iraq. The logic went like this: if G.M. made more
fuel-efficient cars, the country would not need so much oil, and if the country did not need oil,
United States troops would never have invaded.

``I didn't say anything,'' recalled the presenter, who like many others here declined to give her
name because she is not supposed to speak to the news media. ``What can you possibly say? `Thanks?'
''

Even if they ignore the snide comments and occasional jeers, presenting for an ailing car company
just is not as fun as working for one that is thriving. The G.M. and Chrysler spaces are smaller and
less flashy than they were a year ago.

The Jeep exhibit used to have a 54-foot-wide waterfall that continuously dropped 1,000 gallons of
water and was programmed, like an ink jet printer, to spell out brand names and logos in the falling
streams. Not any more.

And for the first time, some G.M. presenters are wearing the same outfits they wore last year.

``We used to get a new one every season,'' said Christine Alt Parry, during a break from her duties
beside a black 2010 GMC Terrain, wearing the flower-pattern dark blazer and black slacks she wore a
year ago. ``I think they're trying to save some money.''

Downstairs at the Kia exhibit, meanwhile, it is party time. A D.J.~is mixing oom-chick-oom-chick
club tracks on an Apple laptop, beside huge LED screens that spell out phrases like ``Schwing!'' and
``Kia Sips Gas.'' The men are wearing new Hugo Boss suits, with dark-purple hankies, and the women
are wearing designer dresses bought recently at the Beverly Center in Los Angeles.

On Friday, Subaru handed out flutes of Brut Cuvee Champagne to visiting Finnish car dealers. And
they are preening over at the Hyundai space, where the staff is decked out in new Armani jackets,
Cole Haan sweaters and a few other items picked up at Nordstrom.

``I haven't seen anyone who looks as sharp as we do,'' said Mark Laffrey, the wardrobe consultant.

The exhibit for Ford, a company in better financial shape than its crosstown rivals in Detroit, is
huge and dominated by an atmosphere that could be described as we-didn't-take-your-money festive.
There are slot car races, a magician doing card tricks, and the crew of MTV's ``Pimp My Ride''
upgrading a car in a cordoned-off section called Mustang Alley, which has a spring-break vibe.

``We're the bad boys of the auto show!'' yells a man who calls himself Flames, one of the
ride-pimpers, as he gets to work.

Not that any of these companies are making huge sums of money. But unlike G.M. and Chrysler, they do
not need to project an air of austerity and seriousness.

So is this new image winning over potential buyers? Well, there are people who say they would indeed
buy from either company. But many attendees echoed the sentiments of Mark Lee, who was photographing
his daughter next to the gleaming rims of a 6,000-pound Hummer. ``Absolutely not,'' he said when
asked if he was tempted.

On the other side of the hall, an electrician, Kurt Moore of Pleasant Valley, N.Y., sat in a Toyota
Highlander and explained why he was not going to buy American anytime soon.

``You know how they say, 'Never buy a car made on a Monday or a Friday?' '' he said, getting
comfortable in the passenger seat. ``It's because the people building the cars aren't focused on the
job. Well, how well do you think they're focused today, with all this talk about how they're going
to lose their jobs?''

The G.M. and Chrysler presenters have heard questions like that, and dozens of variations of it, in
the last couple months.

``We get a lot of, 'You're going out of business,' 'You guys are going bankrupt,' '' said Shannon
Melahn, part of the Chrysler presenting team. She shrugged and added, ``We just smile.''

\section{Obama Signals More Active Response to Piracy}

WASHINGTON -- President Obama vowed Monday to ``halt the rise of piracy'' off the coast of Africa
following the dramatic rescue of an American merchant captain, foreshadowing a longer and
potentially more treacherous struggle ahead as he weighs a series of problematic options.

In permitting members of the Navy Seals to shoot the pirates holding the captain, Richard Phillips,
Mr.~Obama navigated a crisis that played out in full view of the world. But policy makers and
experts said the precision killing of three Somali pirates with three bullets would certainly prove
easier than wiping out the larger threat in the shipping lanes or reversing the instability that
makes Somalia a breeding ground for pirates and Islamic terrorists.

With the ``Black Hawk Down'' episode in Somalia still etched in the American consciousness 16 years
after two helicopters were shot down and 18 American soldiers were killed, Mr.~Obama and his
advisers are wary of becoming deeply involved in the region again. That wariness comes at a time
when Mr.~Obama is already trying to end a war in Iraq and win another in Afghanistan. White House
officials on Monday played down suggestions that the United States could attack pirate bases on
shore, portraying that as premature at best.

Other options that the administration has before it, according to experts, are deploying more ships
to patrol the region, pressing commercial shipping companies to stop paying ransoms and to do more
to defend their vessels, get other nations to help capture pirates and bring them to justice, and
doing more to build up a fledgling transitional government in Somalia.

``All I can tell you is I am confident we will be spending a lot of time in the Situation Room over
the next few weeks trying to figure out what in the world to do about this problem,'' Defense
Secretary Robert M.~Gates told students and faculty members at the Marine Corps War College in
Quantico, Va., according to a military news service. While the national security team discusses the
long-term challenge, the Justice Department said it was deliberating whether to try the lone
surviving pirate in the United States or to turn him over to Kenya for trial. Mr.~Gates said the
four pirates were 17 to 19 years old, but he did not give the age of the one who surrendered to Navy
sailors before the fatal shootout Sunday.

The rescue of Captain Phillips drew widespread praise for the Navy and Mr.~Obama, but some experts
warned that it could escalate the campaign by Somali pirates, who have vowed to take revenge on
Americans and are holding more than 200 hostages from other countries.

Mr.~Obama praised Captain Phillips for his ``courage and leadership and selfless concern for his
crew,'' and he said he was ``very proud'' of the Navy and other American agencies involved in the
operation.

``I want to be very clear that we are resolved to halt the rise of piracy in that region,''
Mr.~Obama said. ``And to achieve that goal, we're going to have to continue to work with our
partners to prevent future attacks. We have to continue to be prepared to confront them when they
arise. And we have to ensure that those who commit acts of piracy are held accountable for their
crimes.''

Senator Russ Feingold, a Wisconsin Democrat and chairman of the Senate's subcommittee on African
affairs, sent a letter to Mr.~Obama calling for a comprehensive strategy to shore up the new
transitional government of Sheik Sharif Sheik Ahmed in Somalia. Mr.~Feingold called central
authority the ultimate solution to piracy.

``People are talking about this as a piracy issue,'' Mr.~Feingold said in a telephone interview.
``That is not the core issue here. It is a symptom of a disunified government.''

Somalia has veered between lawlessness and rule by warlords for most of the past two decades,
providing a haven not only for pirates but also for Islamic militants.

Representative Christopher H.~Smith of New Jersey, the ranking Republican on the House Foreign
Affairs Committee's Africa subcommittee, said any military response would have to be fully supported
and sustained.

``We cannot do what we did in Somalia before,'' he said. ``Brave soldiers died. So everyone,
especially this administration, needs to be careful before pushing that button.''

Bobby Pittman Jr., who was President George W.~Bush's senior adviser on Africa, said that finding
pirates on the open seas would be akin to hunting for ``a needle in the haystack,'' and attacking
camps on the ground would lead to civilian casualties.

``There is no purely military solution to it,'' Mr.~Gates said at the Marine Corps War College.
``And as long as you've got this incredible number of poor people and the risks are relatively
small, there's really no way in my view to control it unless you get something on land that begins
to change the equation for these kids.''

\section{Obama Sees More Pain Now but Hope Later on Economy}

WASHINGTON -- President Obama said on Tuesday that the battered economy was showing signs of
recovery, but he warned Americans that more pain lies ahead and urged them to help build a
foundation for a new, 21st century prosperity.

Speaking just after a disappointing report on March retail sales made it clear that a sustained
recovery is not yet at hand, the president delivered a speech that was part pep talk and part
rebuke, not only for the once high-rolling members of the financial world but for politicians who he
said had deferred tough decisions for too long. ``I want every American to know that each action we
take and each policy we pursue is driven by a larger vision of America's future,'' Mr.~Obama said in
remarks at Georgetown University.

The White House had previewed the event as a ``major speech'' on the economy, but Mr.~Obama did not
break new ground. He did, however, use the occasion to reaffirm his determination to do something
about the rising cost of health care and, later, to shore up Social Security.

The president seemed to guard against being tagged as a ``liberal.'' For instance, he defended his
administration's decision not to take over failing banks: ``Governments should practice the same
principle as doctors: first, do no harm.'' And at another point, he invoked religious imagery.

The president envisioned ``a future where sustained economic growth creates good jobs and rising
incomes; a future where prosperity is fueled not by excessive debt, reckless speculation and fleeing
profit, but is instead built by skilled, productive workers; by sound investments that will spread
opportunity at home and allow this nation to lead the world in the technologies, innovations and
discoveries that will shape the 21st century.''

``That is the future I see. That is the future I know we can have.''

But the near future will bring ``more job loss, more foreclosures and more pain before it ends,''
Mr.~Obama said. Underscoring his point was a Commerce Department report showing that consumer
spending on a wide array of goods declined in March, reflecting a general spirit of uncertainty as
well as continuing job losses.

The president said, as he has repeatedly, that the recently enacted stimulus plan, the efforts to
strengthen the banking system and attempts to rescue the flagging American auto industry have all
borne fruit, demonstrated in part by an increase in home-mortgage refinancings and more lending by
small businesses.

``This is all welcome and encouraging news, but it does not mean that hard times are over,''
Mr.~Obama said, warning that 2009 will be a difficult year, and that no one should expect a return
to full prosperity soon.

As the president spoke, the Federal Reserve Chairman Ben S.~Bernanke told an audience at Morehouse
College in Atlanta that there were ``tentative signs'' that the decline in the economy was slowing.

President Obama called on Americans to take the long view. ``There is no doubt that times are still
tough,'' he said. ``By no means are we out of the woods just yet. But from where we stand, for the
very first time, we are beginning to see glimmers of hope. And beyond that, way off in the distance,
we can see a vision of an America's future that is far different than our troubled economic past.''

Realizing that vision will require a new regulatory structure, one based on 21st century needs
rather than an outdated financial buccaneer ethic, the president said. It will also require work on
deep, complicated issues like health care and energy, he said.

As Mr.~Obama spoke inside Georgetown's Gaston Hall, a small group of abortion opponents demonstrated
outside against the presence at a Catholic university of a president who supports abortion rights.
Using a bullhorn, the protesters could sometimes be heard faintly in the back of the hall during the
president's speech. (There have been similar protests at Notre Dame, where the president is to speak
at commencement exercises on May 17.) Mr.~Obama, alluding to a parable at the end of the Sermon on
the Mount, said he saw a new America whose foundations are built not on sand but on rock, ``proud,
sturdy and unwavering in the face of the greatest storm.''

``We will not finish it in one year or even many,'' he said, ``but if we use this moment to lay that
new foundation, if we come together and begin the hard work of rebuilding, if we persist and
persevere against the disappointments and setbacks that will surely lie ahead, then I have no doubt
that this house will stand and the dream of our founders will live on in our time.''

\section{Faint Signs of Uptick in China Trade}

GUANGZHOU -- Chinese manufacturers of everything from tricycles to truck winches said Wednesday at
the opening of Asia's biggest trade fair that their orders were starting to recover from a steep
plunge over the winter but that demand from Europe and particularly the United States remained
anemic.

That anecdotal evidence of some uptick in trade came as a senior Chinese government economic
forecaster, Fan Jianping, said the economy had grown at an annual pace of ``around 6 percent'' in
the first quarter. That would be the slowest pace since quarterly records began in 1992, indicating
that strong consumer spending and heavy government investments had failed to offset a nose dive in
exports.

The figure given by Mr.~Fan -- the official statistic is to be released Thursday -- was roughly in
line with economists' expectations. It means that the Chinese economy grew much faster over the
winter than any other major world economy, although more slowly than many in China have come to
expect and slower than the 6.8 percent pace in the fourth quarter.

China's slowdown has been felt most keenly by export-oriented factories and their workers, who are
heavily concentrated in the Shanghai area and here in Guangzhou, in southeastern China. By contrast,
``domestic consumption has been holding up very well,'' with auto sales showing particular strength,
said Frank Gong, an economist in the Hong Kong offices of J.P. Morgan.

With Chinese banks pumping out loans at a record pace to stimulate domestic investment spending and
consumption, and with the government spending heavily on new railroads and other projects, the big
question hanging over the Chinese economy is when overseas demand for Chinese goods will recover.

As buyers and suppliers gathered here from around the world Wednesday for the opening of the Canton
Fair, almost no one said that sales had recovered to their levels of a year ago.

Instead, retailers who drew down their inventories to almost nothing over the winter are reluctantly
buying the minimum quantities of goods needed to keep shelves from becoming bare.

And while exporters to Africa, the Middle East and Southeast Asia described clear rebounds in orders
compared with the depths plumbed in December and January, Chinese exporters to Europe and
particularly the United States were much more glum.

Yet compared with the last session of the Canton Fair in late October, there was a clear sense
Wednesday that the worst fears of a collapse in global trade had not been borne out.

``October was the month when we could not see a future,'' said Lin Minli, the sales manager of a
Guangzhou trading company that sells tricycles and other bulky children's toys mainly to South
America and the Middle East.

Mr.~Lin's sales are still down more than 30 percent from a year ago, but that is a big improvement
from a 75 percent drop in October. The improvement ``is because the inventories of our clients have
been depleted,'' he said, adding that ``with the recent rise in raw material prices, some customers
also started putting in more orders because they were afraid prices would rise.''

Chris Sun, a sales manager of Hangzhou Xinweiye Tools, a maker of saw bits in Hangzhou, said that
while exports were still down 30 percent, that was an improvement from the winter. Orders in the
last several weeks suggest further recovery ahead, she added.

``February and March were not good. April does look good; I think the world economy is getting
better,'' Ms.~Sun said, adding that ``the U.S.~market is not so good, but Europe is better.''

Export-dependent countries in Asia in particular have suffered heavily from the global economic
downturn. Many Western retailers halted orders over the winter and drew down their stockpiles
instead; even exporters who did have orders often ran into difficulty finding trade finance because
of the global credit crisis.

Both of those problems now appear to have eased, so the inventory-led plunge in exports appears to
be abating. China announced last week that its exports had fallen 17.1 percent in March from a year
earlier -- a steep decline by most standards, but better than the drop of 25.7 percent in February.

But Dong Tao, an economist in the Hong Kong office of Credit Suisse, said it was much less clear
whether consumer demand for China's goods had really bounced back. ``Whether this pickup in demand
is sustainable, I'm a little more skeptical,'' he said.

In interviews with managers from more than a dozen Chinese companies in a wide range of industries,
clear regional patterns of demand were apparent.

Industrialized countries have tended to lead the world out of previous recessions. But in a reversal
of the usual north-south split, it seems as though emerging markets, often dependent on commodities
exports to China, are faring better this time.

Jet Power International, which makes paper-towel dispensers and other bathroom products in Dongguan,
China, is faring poorly because it relies on the American market. Sales are still falling as
American building managers buy the minimum number of products they need for replacements while new
construction remains weak, said Alex Yu, the company's marketing manager.

Jet Power has just laid off 100 of its 300 workers, and another 50 have been placed on an unpaid
furlough that could last until June, Mr.~Yu said, adding that he was skeptical that other exporters
to the United States were doing much better. ``Those optimistic ones, they are telling lies,'' he
said.

Pessimism among exporters to the United States is consistent with a report from the Commerce
Department in Washington on Tuesday that retail sales fell 1.1 percent in March. Hotel room rates in
Guangzhou are down 40 percent compared with the spring session of the Canton Fair last year, but
rates for the four- and five-star rooms occupied mainly by Americans and other Westerners here are
down even more on very weak demand.

Exporters to developing countries were more upbeat Wednesday. Fujian Yuandong Electric Motor Group
in Fuan, China, has seen a recovery in demand for its electric pumps in Africa, Eastern Europe and
Southeast Asia in the past month, said Orchid Chen, the company's sales director. Sales had been
down 40 percent over the winter compared with a year earlier but are now down 20 percent, she said.

Chinese exporters may be benefiting to some extent from a global version of the ``Wal-Mart effect,''
in which customers turn to lower-cost vendors during times of economic distress.

\section{For Young Japanese, It's Back to the Farm}

YOKOSHIBAHIKARI, Japan -- A motley group of unlikely farmers descended on the countryside here one
recent Sunday, fresh towels around their necks, shiny boots on their feet.

``This is harder than it looks,'' said Tatsunori Kobayashi, a spiky-haired janitor from Tokyo Disney
Resort, as he tromped through a mustard spinach patch with a seed planter, irregular furrows
stretching out behind him.

He is part of Japan's 2,400-strong Rural Labor Squad, urban trainees dispatched to the countryside
under a pilot program to put Japan's underemployed youth to work tilling its farms.

Started last month as part of Prime Minister Taro Aso's stimulus plans, the program stems from
growing concern about both the plight of Japan's younger workers and the dismal state of farms. In a
play on words, the squad's name in Japanese -- Inaka-de-hatarakitai -- is also its rallying cry:
``We want to work in the countryside!''

The predicament of Japanese in their 20s and 30s dates back to the lost decade of the 1990s, when
many failed to find good, stable work. Today, a disproportionate number endure low-wage jobs -- a
potential portent for America's students and first-time job seekers plunging into a shallow job
market in the United States.

As the Japanese recession has worsened, younger workers have taken the brunt of wage cuts and
layoffs, especially in manufacturing. Now the government views the slump -- Japanese exports fell
almost 50 percent year-to-year in February -- as a chance to divert idle labor to sectors that have
long suffered from worker shortages, like agriculture. Many young Japanese, for their part, have
shown a growing interest in farming as disillusionment rises over the grind of city jobs and
layoffs. Agricultural job fairs have been swamped with hundreds of applicants; one in Osaka
attracted 1,400 people.

``Young people want jobs, and farmers need the extra hands,'' said Isao Muneta, an agriculture
ministry official who coordinates the 1.3 billion yen (\$13 million) program, part of a larger
stimulus package. ``It's the perfect match.''

Whether it will save Japan's deteriorating economy is something else. ``Rural communities could
benefit from an influx of young people,'' said Masashi Umemoto at the National Agricultural Research
Center. ``But it's unrealistic to look to agriculture as a solution to the country's unemployment
problems.''

He added, ``There aren't enough farming jobs.''

Like the French and the British, whose industrial societies have deep (if distant) rural roots, the
Japanese have long romanticized life in the countryside. Only 4 percent of Japan's labor force works
in agriculture, but a reverence for the country's rice-farming heritage is strong. Japanese children
grow up with warnings not to waste a single grain of rice, out of respect for farmers' labor. In an
annual ritual, the Japanese emperor makes an offering of rice harvested from paddies within the
palace grounds to Shinto deities. And in international trade talks, rice remains the most sensitive
crop for Japan.

Beneath this romanticism, however, is a stark reality. Japanese farming is a picture of
inefficiency, and the rural work force is graying. A decline in rice prices has hit farms hard --
only the largest farms still turn a profit from harvesting rice, forcing farmers to take on extra
jobs. The farms most desperate for workers do not have the means to pay for new recruits.
Agricultural jobs pay as little as \$1,500 a month and are often seasonal.

Overgrown plots abound in Yokoshibahikari, a town of 26,000 about 43 miles east of Tokyo.

``We're all old folk and thankful to have young people come help us,'' said Hitoshi Suzuki, 57, and
head of a cooperative of family farms that share equipment to reduce overhead costs. (One of the
cooperative's farmers is 83.)

Rural communities themselves effectively shut out new blood by making it difficult for outsiders to
set up their own farms, says Takayuki Yoshioka, a coordinator at the nonprofit organization that
runs the Yokoshibahikari program. People with no local links who want to buy farmland are subjected
to a vetting process by local farming committees that can take years.

``I believe the possibilities are limitless in agriculture,'' said Mr.~Yoshioka, who is interested
in starting his own farm. ``But there are also big barriers.''

Shinji Akimoto, who until recently worked in information technology, is not intimidated.

Fearful of constant staff cuts as business deteriorated, Mr.~Akimoto, 31, quit his job last month
and days later started training in Yokoshibahikari. His three-day, government-financed training
program has been a succession of whirlwind lessons in rice and vegetable planting, cleaning pig
sties and feeding cattle.

``I had nothing much to lose, and in times like these, I felt I needed to learn to make my own
living,'' he said. He chuckled and twirled a finger in the air. ``Did you know pigs really do have
curly tails?''

Mr.~Akimoto's team of 10 is a hodgepodge: the Disney janitor, a recently laid-off landscape artist
and several college students. They all get 7,000 yen a day, about \$70, and free food and board.

They all shared a common complaint: there was no convenience store nearby for drinks and snacks. One
trainee persuaded a farmer to lend him his light truck, so he could get cigarettes.

``My friends think I'm crazy for coming here,'' said Tomoka Inoue, 20, a management major who said
she was widening her job search to include farming. ``But I think people are becoming more aware of
where our food comes from, and I want to get more involved with that.''

Experts say the program's wider economic impact will be limited in the face of the severe challenges
facing Japan's economy: gross domestic product shrank at an annualized rate of 12.1 percent in the
fourth quarter of 2008, and unemployment is at a three-year high of 4.4 percent.

But the government is going ahead with plans to begin yearlong farm intern placements later this
year. Increasing agricultural employment is part of a new \$154 billion stimulus package that
Mr.~Aso announced last week.

Mr.~Kobayashi, the janitor at Disney, says his time as a trainee has helped him decide he wants to
take up farming leeks, this town's main crop. He intends to take another week off to train with a
local leek farmer, Yoshinori Yamazaki, who is looking for someone to take over his farm.

``This is just too perfect,'' Mr.~Kobayashi gushed. He said leeks were his favorite vegetable, and
he had read that they were easy for beginners to grow and bring in a stable income.

But Mr.~Yamazaki, the leek farmer, was skeptical. ``You can't learn farming in just a year, or even
several years. It's a lifetime profession,'' he said. ``I worry this is just a fad. I'm worried that
when the economy picks up, they'll all flock back to the city.''

\section{The Quest To Shrink The S.L.R.}

Physics can be such a bummer, you know? Those doggone rules of nature make our luggage heavy, melt
polar ice caps and conspire to make the toast land jelly side down.

And physics explains why you can't take professional-quality pictures with a small camera. Big,
bright, beautiful photos require big sensors and big lenses to shine on the chips. So we're stuck
with two camera categories: pocket ones that take mediocre pictures and big heavy black ones
(S.L.R.'s) that take stunning photos but require a neckstrap or, ideally, a wheelbarrow.

Because both ends of that spectrum are already teeming with camera models, the manufacturers are
eager to differentiate in hopes of finding new markets. This month, Nikon and Canon are going for a
category that, so far, is sparsely populated -- the one right in the middle.

Call it the mini-S.L.R.. Nikon's new D5000 (\$850 with lens) and Canon's PowerShot SX1 IS (\$600)
look very similar. The bodies are sized and shaped identically -- like mini-S.L.R.'s --and they
weigh roughly the same (around 1.5 pounds).

Each camera can capture high-definition video; each even has a mini-HDMI jack for connecting to your
hi-def TV. (Mini-HDMI requires a special cable, not included.)

Each has a flip-out, swiveling screen on the back. That feature permits you to shoot at high, low
and side angles, not to mention allowing self-portraits. (Nikon's screen flips downward below the
camera, rather than off to the side -- an arguably better design, since it keeps your view in line
with the lens.) But a flip-out screen also gets the big black camera away from your face, thereby
getting more natural, relaxed expressions from your subjects, especially young ones.

Despite these similarities, though, these cameras approach the middle from opposite ends. Nikon has
tried to consumerize an existing S.L.R.; Canon has tried to S.L.R.-ize an existing consumer camera.

The Nikon D5000 is a genuine S.L.R. It takes interchangeable lenses, its eyepiece lets you peer out
through the lens itself, it has an enormous light sensor inside, and it has no shutter lag (the
delay after you press the button).

In fact, if you just read the brochure, you might mistake the D5000 for Nikon's more advanced, more
expensive D90 (\$1,150 with lens); the sensor inside, the specs and most of the features are exactly
the same.

The D5000 has the same superb 12.3-megapixel sensor, the same sensor-shaking mechanism to eliminate
dust spots, and the same 11-point autofocus system, which makes it easier to compose focused shots.
It even has the same burst rate (4.5 shots a second), which is fantastic not just for sports and
action but for capturing fleeting facial expressions.

Like most S.L.R.'s these days, the D5000 also has Live View mode. It makes focusing slower, but it
means that you can frame your shots using the screen, if you like, instead of holding the camera to
your eye. The D90 was the first video S.L.R.; the D5000 inherits that hi-def goodie, too.

But here's where things get crazy: Nikon has packaged all of these professional features into a much
smaller, lighter camera that's obviously aimed at amateurs. For example, the new camera is loaded
with scene modes to make the thing less intimidating to novices -- canned presets for Sunset, Pets
or Beach/Snow, each with a sample photo on the screen.

As on other recent low-end Nikon models, the D5000 also makes photography easier to understand by
displaying a visual representation of the actual aperture (lens opening) inside as it grows and
shrinks. The screen also indicates your shutter speed both graphically and with a fraction --
``1/25,'' for example, rather than the ``25'' that appears on most S.L.R. screens. (The D5000 lacks
a top-panel status screen like the big boys have -- one casualty of the smaller body.)

The D5000, in other words, is such a mix of high- and low-end that Nikon's representatives actually
asked my opinion on how to describe the thing. Here's my suggestion: it's a professional camera
trapped in a consumer-cam's body.

A few perks of the D90 are missing on the D5000. You lose the ability to use some older Nikon
lenses. The battery is smaller, good for 510 shots on a charge instead of 850. The camera doesn't
communicate with wireless flash attachments.

There are some consolations, though. A new Quiet Mode slightly muffles the clack of the shutter.
Focusing before you shoot a movie is faster, although, tragically, you still can't change focus
while filming a movie.

The upshot here is that the D5000 produces spectacular, sharp, fast, brilliant photos in all kinds
of light. It has the same expensive guts of Nikon's higher-end models in a small, rugged, lovingly
designed, user-friendly body -- and it's a home run.

You can't say the same for Canon's new PowerShot SX1. It's an evolution of the PowerShot S IS series
that arrived in 2004; the IS stands for image stabilized, and the S indicates a camera that's
equally good at movies and stills.

Unlike most cameras, for example, the SX1 can change focus and even zoom while you're filming -- the
lens motors are so quiet, the excellent stereo microphones don't pick up its sound. It even has
separate shutter buttons for video and photos, so you can snap a still in the middle of recording
video.

What made the SX1 such a hotly anticipated model, though, is its light sensor. It's a CMOS chip, the
type that's found in S.L.R.'s, including Canon's popular Rebel series. It's much more powerful than
the chip found in regular consumer cameras. That doesn't do anything for photo quality, but it gives
the SX1 several features that its predecessors lacked, like a very high burst mode (4 frames a
second) and full high-definition movies.

All that and a wide-angle, 20X optical zoom lens (that's not a typo). Where do we sign?

Unfortunately, the SX1 simply doesn't take great pictures. It's a feature list in search of a
camera. In bright sunlight, you'll adore the thing. It's fast, it's easy, it's always ready, and the
photos are sharp and saturated. But as the light dims, so does the quality.

Noise (graininess) begins to appear even in outdoor overcast shots (say, ISO 400) -- and when the
camera's software tries to iron out that graininess, you get soft edges and lost detail.

Unless you dial down the exposure, the SX1 also tends to blow out the highlights -- that is, it
turns brighter areas of your photos into bleached blobs of nuclear white. (You can see samples from
both of these cameras at nytimes.com/tech.)

The SX1 also has, in place of a glass eyepiece, an electronic viewfinder (a tiny screen inside the
eyepiece). That's O.K. if that little screen has excellent resolution, but this one is so coarse,
you actually see the pixels. Not good.

Camera-testing Web sites like dpreview.com and testfreaks.com have compared the SX1 with its
less-costly sibling, the SX10 (\$340). The SX10 lacks the CMOS chip and hi-def video -- but takes
better photos. Weird, huh?

It's so unlike Canon, which makes the world's best small cameras, to introduce a model with these
compromises. It's so cool to have a 20X zoom available, a dedicated button for switching between
widescreen and standard modes, all those manual controls and the option to use standard AA batteries
in a pinch. But for \$600, you deserve a camera with a sensor that is up to the light-gathering
task.

What the world really wants, of course, is a whopping big sensor like the D5000's -- in a
shirt-pocket camera. Unfortunately, you probably won't find a chip like that in a camera much
smaller than the D5000 itself -- at least not until there's an amendment to the laws of physics.

\section{In Recruiting an Afghan Militia, U.S.~Faces a Test}

MAIDAN SHAHR, Afghanistan -- The ambitious American plan to arm local militias in villages across
the country was coming down to a single moment.

The American officers sat on one side of a long wooden table; a group of Afghan elders on the other.
The pilot program was up and running, but the area's big enclave of Pashtuns -- the ethnic group
most closely identified with the Taliban -- had not sent any volunteers. The Pashtuns were worried
about Taliban reprisals.

``We agreed to meet today and, I believe, make a decision,'' Lt. Col. Kimo Gallahue told the 11
elders. ``Time is running out.''

Then he laid down a challenge: ``I am so proud to be in the same room with the men who defeated the
Soviet Union. Please find the courage to take responsibility for your own security.''

The elders, in their turbans and beards, stared blankly at the Americans across the table.

For two hours, the meeting unfolded, laying bare the torments facing any Afghan Pashtuns who might
be contemplating defying the Taliban -- and the extraordinary difficulties facing American officers
as they try to reverse the course of the war.

The meeting in Maidan Shahr, Wardak Province's capital, tucked into the mountains about 30 miles
southwest of Kabul, concerned one of the most unorthodox projects the Americans have undertaken here
since the war began in 2001: to arm, with minimal training, groups of Afghan men to guard their own
neighborhoods.

The military is borrowing a page from a similar program that helped bring about the recent calm to
Iraq, where the Americans signed up more than 100,000 Iraqis, most of them Sunnis and many of them
insurgents, to keep the peace.

The hope here is that the militias will come to the aid of the overwhelmed Afghan Army and the
police, which take longer to train and equip and number only about 160,000. Hundreds were killed
last year in Taliban attacks.

If the militias work in Wardak, the Americans say they want to replicate them throughout the
country. So the experience in Wardak has been instructive, for what the Americans can accomplish and
what they cannot.

At first, everything went well in Jalrez, the mountainous area where the program is based. Young men
from two of Jalrez's main ethnic groups, the Tajiks and the Hazaras, enthusiastically came forward;
both have largely supported the American presence. Several dozen Pashtuns from other villages showed
up as well. Two hundred forty-three volunteers were selected, each vetted by the police, the elders
and the local religious leaders. The first crop of recruits went through the three-week course --
presided over by American Special Forces officers -- and graduated three weeks ago. They are now
patrolling the dirt roads of Jalrez.

The trouble came from the Pashtun enclave of Zayawalat, one of five large villages in Jalrez.

The Americans setting up the guard force waited patiently, hoping to bring Zayawalat's elders along.
They agreed to a meeting with the elders, and then another and another. At a meeting last week, the
fourth, the Pashtun elders said they would make a final decision and report back this week.

But when they showed up Monday morning, the elders said they still were not ready to give up their
sons. ``It's not that the people in Zayawalat don't support the government -- they do,'' said Hajii
Janan, the leader of the Wardak provincial council, who presided over the meeting. ``But, as you can
see, people are under pressure.''

Mr.~Janan was not exaggerating. Last month, a local Taliban commander, Abdul Jameel, based in Maidan
Shahr, came forward with 10 of his fighters and declared that he would fight no more. Wardak's
governor, Halim Fidai, accepted his surrender and told him to go home. The governor offered
Mr.~Jameel no protection for this act of defiance of the Taliban. Two weeks ago, Taliban gunmen
entered Mr.~Jameel's home and killed him, his wife, his uncle, his brother and his daughter.

At the meeting on Monday, the elders produced copies of three different leaflets, called ``night
letters,'' that had been slipped under the doors of their homes. ``The holy warriors urge you to
reject the American plan,'' one of the milder leaflets said. ``Stay with us.''

The Americans said that although they were sympathetic to the Pashtuns' fears, the time for bravery
had come. In January, the Americans dispatched two battalions, about 1,600 men, to Wardak Province,
a huge increase over what was here before. Afghans had to risk their lives, too.

The Americans' staunchest ally was the local commander of the National Directorate of Security, Gen.
Mullah Razik. He is a Pashtun, too, but has been working with the Americans since 2001.

``This isn't some program imposed by the Americans,'' General Razik told the elders. ``It is for
you, your country, your sons and your daughters. It's up to us.''

Then there were the two leaders of the council of elders, Ghulam Mohammed and Abdul Ahmed. Both had
fought for the Taliban. They were leading the opposition to the militia plan.

Like most Afghans, the life stories of Mr.~Mohammed and Mr.~Ahmed are full of twists and turns and
turnabouts, making their current allegiance somewhat difficult to determine. In the 1980s,
Mr.~Mohammed and Mr.~Ahmed joined a group, the Movement for the Islamic Revolution, and fought the
Soviet Union. One of their comrades then was General Razik. All three men have known one another
since they were young.

Then, in the 1990s, the Taliban came along, and Mr.~Mohammed and Mr.~Ahmed joined them. ``The
Taliban were forcing people to join,'' Mr.~Mohammed said before the meeting. Mr.~Mohammed and
Mr.~Ahmed said they quit the Taliban in November 2001, when the group was driven from Kabul by the
American-led invasion.

But one of the questions hanging over the meeting was whether Mr.~Ahmed and Mr.~Mohammed were still
working with the Taliban, after all.

``I don't think they even know what side they are on,'' said an American Special Forces officer, who
is helping to train the militias. He spoke on the condition that he not be named. ``Those two are
right on the fence, trying to figure out which side to go with. They are trying to get the best deal
for themselves.''

Mr.~Ahmed said as much, wondering whether, if he threw in his lot with the Americans and the Afghan
government, they could keep him safe.

``Sure, I could round up 50 guys for you,'' Mr.~Ahmed said. ``But the general here has a uniform.
He's got guns. All I've got is the turban on my head and a piece of bread.''

After two hours, the elders could not decide. Once again, they asked the Americans for more time.

The Americans shook their heads in exasperation.

``This is your last chance,'' General Razik told the elders. ``If you don't take it, we are just
going to associate you with the Taliban.''

And they agreed to meet again.



\section{No Easy Workout}

FAYETTEVILLE, Ark. -- Rick Williamson, a Chicago banker turned junk-loan buyer, knew the
name-calling would start again when he moved to foreclose on the Fayetteville Athletic Club, a
sprawling, family-run gym and spa complex in this corner of northwest Arkansas.

``Vulture, bottom feeder,'' Mr.~Williamson said, recalling insults thrown his way during his years
hunting for bargains in distressed real estate.

And just as he predicted, the insults are flying.

``Trash-eating rat,'' is how Robert Shoulders, the club's owner for the last 13 years, described
Mr.~Williamson. ``What he does is reprehensible. I am not sure how he can sleep at night.''

Mr.~Williamson sees it differently. And the government agrees.

As he points out, it isn't his fault that Mr.~Shoulders overextended himself by borrowing \$10
million to add the spa, a preschool, tennis courts and other amenities to his once-modest health
club. And by moving aggressively to collect on loans like these, Mr.~Williamson says, he is playing
a crucial role in helping clean up the bad debts that are clogging the economy.

``If you want to fix what is ailing this country, you need to destroy the worthless debt out
there,'' he said, adding with characteristic candor, ``You have to shut down the zombies, liquidate
their assets.''

Clashes like this are increasingly taking place across the country, as banks struggle through their
worst crisis in a generation. And Mr.~Williamson is part of a niche industry that buys at
bargain-basement prices the hard-to-collect commercial loans that the Federal Deposit Insurance
Corporation auctions off after it seizes a failed bank. They take on the most troubled assets, the
ones others aren't willing to touch, in exchange for potentially lucrative returns.

Nationwide, the F.D.I.C. has seized 58 banks over the last 15 months and through loan auctions has
sold about \$2.2 billion worth of loans. The \$1.2 billion generated for the government has gone
into the fund to insure deposits at other banks.

Many of the auctioned loans are considered ``performing,'' meaning the borrowers are current on
their payments. But a portion of them are delinquent, like the one held by Mr.~Shoulders and his
wife, who stopped making payments after their local bank, ANB Financial of Bentonville, Ark., failed
last year.

Loath to throw homeowners out on the street, the F.D.I.C. has made it a national policy to try to
avoid foreclosures on the single-family mortgages it inherits from failed banks, even if payments
are past due. But there are no such rules to protect struggling small businesses, whose loans are
more typically sold at the F.D.I.C. auctions.

The borrowers, like the Shoulderses, suddenly find themselves dealing with aggressive out-of-state
loan consolidators, rather than local banks with long-standing community ties. And the
consolidators' primary objective often is to immediately collect on the debt, or seize the
collateral.

The single biggest buyer at these auctions has been Andrew Beal, a banking billionaire from Texas
who made his fortune buying distressed debt during the savings and loan crisis.

By the end of February, Mr.~Beal had paid more than \$200 million to buy \$438 million worth of
loans, according to agency records. Some of the loans came from the failed Arkansas bank.

Mr.~Beal is hardly averse to risk. He is famous for trying (and failing) to build his own space
satellite launch company, and for luring some of the world's best poker players to a series of
games, with him as a participant, and betting pots worth \$2 million. Mr.~Beal, in a telephone
interview, said he went to great lengths not to push people out of their businesses, but at times he
had no choice.

``Borrowers force us into litigation,'' he said. ``They don't want to perform on their loan, they
won't talk to our workout people. What are we supposed to do, send them a vase of roses?''

Another buyer here, Daniel C.~Cadle of Newton Falls, Ohio, has drawn some of the loudest protests
over the years for hardball collection practices.

In 2007, Massachusetts authorities, in denying him a license to operate in the state, accused
Mr.~Cadle of not having ``the character, reputation, integrity and general fitness to engage in the
business of a debt collector in an honest, fair and sound manner.'' In one instance they cited, a
representative of Mr.~Cadle's firm told a debtor that he had photographs of the man outside his
house, knew his wife's name and had information on his bank account.

For nearly a decade, Mr.~Cadle faced the threat of arrest by authorities in Texas as the result of a
lengthy debt collection dispute in which a county court concluded that Mr.~Cadle had tried to
financially ruin a debtor and used a pattern of abusive lawsuits.

Mr.~Cadle says his run-ins with the law are a consequence of his willingness to demand repayment
from some powerful people. And he makes no apologies for working hard to collect from debtors.

``We are heartless because people refuse to pay us even payments they can afford?'' he said. ``You
aren't taking advantage of anyone if you are simply asking them to pay what they owe you.''

Some of the debts are so hard to collect that the F.D.I.C. lets these loans go for a song. LeMire
Schmeglar, a mortgage broker in Chicago, bought 191 delinquent loans with a book value of \$6
million. He paid just over \$15,000.

F.D.I.C. officials, in a written statement to The New York Times, said the agency had tried to
balance its mandate to recoup as much as possible through the sale of loans it acquired while still
protecting the interests of borrowers.

``We are statutorily required to sell assets in a way that will get as much money back for the
uninsured depositors and other creditors, including our insurance fund,'' Andrew Gray, an F.D.I.C.
spokesman, said in the statement. ``By focusing on a small percentage of particularly distressed
assets, the view becomes distorted.'' Mr.~Gray added that the agency was creating a new unit to
field any complaints from customers who had loans at failed banks.

There is a fair amount of mystery about the F.D.I.C.'s auction process. Bidding is not open to the
public. And the list of winners is next to impossible to decipher. It is filled with enigmatic names
like ``Brown Bark III'' and ``Oceanside, CA 92054,'' entities that have no Web sites or telephone
listings.

Matthew Anderson, a partner at Foresight Analytics, a California firm that has studied the F.D.I.C.
sales, said the buyers of these loans typically hoped to earn a profit of 25 to 35 percent a year --
a high rate of return, but achievable only by buying loans with big risks. The forceful collection
efforts are spreading nationally as bank failures increase. A search of court records shows cases
filed by buyers of F.D.I.C. business loans in states like Florida, Arizona and Michigan.

In Arizona, Michael W.~Bauer, 53, an electrician and owner of a vending machine business, said that
after 15 years in business, and never once defaulting on a loan, he might have to file for
bankruptcy because an entity called SMS Financial XVII, which bought his loans from the F.D.I.C.
after the failure of First National Bank of Nevada, had filed a foreclosure lawsuit against him.

``I owe the money. I know that, and I am prepared to repay my loan over time,'' he said. ``But to be
strong-armed like this, at the moment when you are the most vulnerable, it just does not seem
right.''

Northwest Arkansas, home to Wal-Mart and the poultry giant Tyson Foods, is hardly the epicenter of
the banking crisis. But because ANB Financial was one of the first reasonably large banks to go
bust, just as the wave of mortgage defaults was getting under way, the distressed loan drama is
playing out here with particular intensity.

Even before the F.D.I.C.- sponsored auctions took place last fall, some of the agency's Arkansas
staff tried to warn small-business owners in the Fayetteville area to try to quickly clear up any
past-due debts or be prepared to face the consequences, according to company owners and a former
agency contractor.

``They said it out loud and quite clearly,'' said Marsha G.~Dunbar, who worked as a contractor for
the F.D.I.C. after it took over ANB, where she once was a loan officer. ``Settle up these debts now,
or it is not going to be pleasant,'' Ms.~Dunbar said, paraphrasing the warnings.

Mr.~Shoulders and his wife, Katherine, who bought the Fayetteville Athletic Club 13 years ago, heard
these same warnings. They offered to pay \$6 million immediately, and an additional \$1 million upon
the future sale of the gym, if the agency would agree to forgive their \$10 million in debt.

F.D.I.C. officials, seeking to maximize their return for the insurance fund, said they proposed to
reduce the loans' interest rates instead, an offer Mr.~Shoulders disputes was ever made. So the
loans were bundled with others and auctioned off in late October.

That was where Mr.~Williamson came in.

More than a year ago, at the first signs of recession, Mr.~Williamson, a one-time Illinois soybean
farmer who became a banker in the 1980s, quit his corporate bank job to form his own business
specializing in buying up distressed debt.

Like many in this business, Mr.~Williamson, 57, is not shy about confrontation. He is openly
disdainful of what he calls ``the nanny state,'' referring to a federal government that he says too
often tries to protect delinquent borrowers.

He bought the Shoulderses' loans for 34 cents on the dollar, an outcome that left the gym owners
furious, because the auction produced a price that was far lower than they had offered to pay the
F.D.I.C. From the start, relations between the two sides were tense.

``These guys put a gun to our head and said 100 percent is due in 10 days,'' said Mr.~Shoulders, 54,
a tall, wiry man who uses a giant green exercise ball as his desk chair and wears a big name tag on
his gym fleece that simply says ``Bob.''

Within a matter of weeks, a foreclosure lawsuit was filed, as the Shoulderses, citing credit markets
that had grown even tighter, said they could raise only \$5 million to repay their loans. After
evaluating how much the collateral -- the lavish club -- was worth, Mr.~Williamson rejected the
offer. The lawsuit is pending in the local court system.

Business is still brisk at the Fayetteville Athletic Club. Its workout rooms are filled with local
doctors, lawyers, Wal-Mart executives, even the local tax assessor. But the Shoulderses, who are
prominent figures in the community here, soon may no longer own the club.

``Were we overextended? Yeah, I will not deny it,'' Mr.~Shoulders said, as he walked through the
gym, greeting members by name. ``But we have put our heart and soul, and every penny we have earned,
into building this place. It is criminal the way we have been treated.''

As Mr.~Williamson sees it, a debt is a debt.

``When I grew up, people paid for their mistakes,'' he said.

\section{Turning Tables, U.S.~Troops Ambush Taliban With Swift and Lethal Results}

KORANGAL OUTPOST, Afghanistan -- Only the lead insurgents were disciplined as they walked along the
ridge. They moved carefully, with weapons ready and at least five yards between each man, the
soldiers who surprised them said.

Behind them, a knot of Taliban fighters walked in a denser group, some with rifles slung on their
shoulders -- ``pretty much exactly the way we tell soldiers not to do it,'' said Specialist Robert
Soto, the radio operator for the American patrol.

If these insurgents came close enough, the soldiers knew, the patrol could kill them in a batch.

Fight by fight, the infantryman's war in Afghanistan is often waged on the Taliban's terms.
Insurgents ambush convoys and patrols from high ridges or long ranges and slip away as the
Americans, weighed down by equipment, return fire and call for air and artillery support. Last week
a patrol from the First Infantry Division reversed the routine.

An American platoon surprised an armed Taliban column on a forested ridgeline at night, and killed
at least 13 insurgents, and perhaps many more, with rifles, machine guns, Claymore mines, hand
grenades and a knife.

The one-sided fight, fought on the slopes of the same mountain where a Navy Seal patrol was
surrounded in 2005 and a helicopter with reinforcements was shot down, does not change the war. It
was one of hundreds of firefights that have occurred in the Korangal Valley, an isolated region
where local insurgents and the Americans have been locked in a bitter stalemate for more than three
years.

But as accounts of the fight have spread, the ambush, on Good Friday, has become an emotional
rallying point for soldiers in Kunar Province, who have seen it as a both a validation of their
equipment and training and a welcome bit of score-settling in an area that in recent years has
claimed more American lives than any other.

The patrol, 30 soldiers from the First Battalion, 26th Infantry, had left this outpost before noon
on April 10, and spent much of the day climbing a ridge on the opposite side of the Korangal River,
according to interviews with more than half the participants.

Once the soldiers reached the ridge's crest, almost 6,000 feet above sea level on the side of a peak
called Sautalu Sar, they found fresh footprints on the trails, and parapets of rock from where
Taliban fighters often fire rifles and rocket-propelled grenades down onto this outpost.

The platoon leader, Second Lt. Justin Smith, selected a spot where trails intersected, and the
platoon dug shallow fighting holes before dark. Claymore antipersonnel mines were set among the
trees nearby.

At sunset, Lieutenant Smith called for a period of absolute silence, which lasted into darkness.
Then he ordered three scouts to sit in a listening post about 100 yards away, 10 feet off the trail.

The scouts set in. Less than a half-minute later, a column of Taliban fighters appeared, walking
briskly their way.

Sgt. Zachary R.~Reese, a sniper, whispered into his radio. ``We have eight enemy personnel coming
down on our position really fast,'' he said. He could say no more; the Taliban fighters were a few
feet away.

More appeared. Then more still. The sergeant counted 26 gunmen pass by.

The patrol, Second Platoon of Company B, was in a place where no Americans had spent a night for
years, and it seemed that the Afghans did not expect danger.

The soldiers waited. The rules of the ambush were long ago drilled into them: no one can move, and
no one can fire until the patrol leader gives the order. Then everyone must fire at once.

The third Taliban fighter in the column switched on a flashlight, the soldiers said, and quickly
switched it off. About 50 yards separated the two sides, but Lieutenant Smith did not want to start
shooting too soon, he said, ``because if too many lived then we'd be up there fighting them all
night.''

He let the Taliban column continue on. The soldiers trained their weapons' infrared lasers, which
are visible only with night-vision equipment, on the fighters as they drew closer. The lasers mark
the path a bullet will fly.

The lead fighter had almost reached the platoon when Pvt. First Class Troy Pacini-Harvey, 19, his
laser trained on the lead man's forehead, moved his rifle's selector lever from safe to
semi-automatic. It made a barely audible click. The Taliban fighter froze. He was six feet away.

Lieutenant Smith was new to the platoon. This was his fourth patrol. He was in a situation that
every infantry lieutenant trains for, but almost no infantry lieutenant ever sees. ``Fire,'' he
said, softly into the radio. ``Fire. Fire. Fire.''

The platoon's frontage exploded with noise and flashes of light as soldiers fired. Bullets struck
all of the lead Taliban fighters, the soldiers said. The first Afghans fell where they were hit, not
managing to fire a single shot.

Five Taliban fighters bolted to the soldiers' left, unwittingly running squarely into the path of
machine-gun bullets and the Claymore mines. For a moment, the soldiers heard rustling in the brush.
They detonated their Claymores and threw hand grenades. The rustling stopped.

Two other Taliban fighters had dashed to the right, toward an almost sheer drop. One ran so wildly
in the blackness that his momentum carried him off the cliff, several soldiers said.

Another stopped at the edge. Pvt. First Class Brad Larson, 19, had followed the man with his laser.
``I took him out,'' he said.

The scout at the listening post shot three of the fleeing fighters, and dropped two more with hand
grenades. ``We stopped what we could see,'' Sergeant Reese said.

The shooting had lasted a few minutes. The hillside briefly fell quiet. The surviving Taliban
fighters, some of whom had run back up the trail, began shouting in the darkness. ``We could hear
them calling out to one another,'' Specialist Soto said.

Lieutenant Smith called the listening post back in. After two Apache attack helicopters showed up,
an F-15 dropped a bomb on the Taliban's escape route, about 600 yards up the trail. Then the
lieutenant ordered teams to search the bodies they could find on the crest.

Sergeant Reese gave his rifle to another sniper to cover him while he tried to cut away a Taliban
fighter's ammunition pouches with a four-inch blade. The fighter had only been pretending to be
dead, the soldiers said. He lunged for Sergeant Reese, who stabbed him in the left eye.

In all, the soldiers found eight bodies on the crest. They photographed them to try to identify them
later, and collected their weapons, ammunition, radios and papers. Then the patrol swept down a
gully where a pilot said he saw more insurgents hiding.

Four scouts, using night-vision gear, spotted five fighters crouching behind rocks, and killed them
with rifle and machine-gun fire, the scouts said. The bodies were searched and photographed, too.
The platoon began to hike back to the outpost, carrying the captured equipment.

Second Platoon, Company B has endured one of the most arduous assignments in Afghanistan. Eight of
the platoon's soldiers have been wounded in nine months of fighting in the valley, part of a bitter
contest for control of a small and sparsely populated area.

Three others have been killed.

In a matter of minutes, the ambush changed the experience of the surviving soldiers' tours. The
degree of turnabout surprised even some the soldiers who participated.

``It's the first time most of us have even seen the guys who were shooting at us,'' said Sgt. Thomas
Horvath, 21.

The next day, elders from the valley would ask permission to collect the villages' dead. Company B's
commander, Capt.~James C.~Howell, would grant it.

But already, as the soldiers slid and climbed down the mountain, word of the insurgents' defeat was
traveling through Taliban networks.

Specialist Robert C.~Oxman, 21, had put a dead fighter's phone in his pocket. As the platoon
descended, the phone rang and rang, apparently as other fighters called to find out what had
happened on Sautalu Sar. By sunrise, it had been ringing for hours.

\section{Obama to Push to Ratify Treaty on Gun Trafficking}

MEXICO CITY -- President Obama, vowing to confront drug cartels that are ``sowing chaos in our
communities,'' called Thursday for the Senate to ratify a long-stalled treaty aimed at curbing
illegal arms trafficking. But Mr.~Obama also suggested that he would not press lawmakers to revive
an expired ban on assault weapons.

On the first day of a four-day swing through Latin America, Mr.~Obama arrived here on Thursday
afternoon for a visit that was intended to demonstrate what he called ``a new era'' in United
States-Mexican relations and to show support for his Mexican counterpart, President Felipe Calderón.
Stemming the rising tide of drug violence was the centerpiece of their agenda.

But while the two pledged their ``shared responsibility'' in the drug wars, they seemed to diverge
on the issue of the assault weapons ban, which expired in 2004. Mr.~Calderón said 90 percent of the
weapons seized in Mexico can be traced to the United States, adding that the flow of weapons
increased after the ban expired.

Mr.~Obama campaigned on a platform of renewing it, but he suggested Thursday that reviving the ban
was politically impossible because of opposition from gun enthusiasts. During a news conference with
Mr.~Calderón, the president said he thought there were more immediate ways to curb the flow of arms.

``None of us are under any illusion that reinstating that ban would be easy,'' Mr.~Obama said, while
insisting he was ``not backing off at all'' from his conviction that renewing the ban made sense.
Instead, he said he was taking other steps, including sending more border agents to patrol the
border with Mexico and more helicopters to Mexico. He also said he was pressing the Senate to ratify
the long-stalled inter-American arms treaty, which the United States had a major hand in negotiating
through the Organization of American States.

The treaty, signed by President Bill Clinton, languished in the Senate, but went into effect in 1998
after two dozen other nations ratified it. The treaty seeks to crack down on illicit firearms by,
among other things, establishing a system for the import, export and transfer of firearms, and by
fostering cooperation among law enforcement agencies investigating illegal trafficking.

``It makes good policy sense,'' said Peter DeShazo, a former senior State Department official who
has advocated ratifying the pact. ``It's very hard for the United States to call on other countries
to cooperate on controlling the flow of illegal arms if we haven't ratified a major inter-American
convention.''

But it was unclear what type of reception the treaty would receive now in Washington. Aides to
Mr.~Obama said he was working closely with Senator Harry Reid, the Democratic leader, to get the
treaty ratified, but Mr.~Reid issued a statement that offered no commitment to bring the treaty up
for a vote.

On Capitol Hill, a senior Democratic official said that despite the president's urging, it would be
difficult to move forward on the long-stalled treaty given the Senate's already crowded agenda, as
well as a continuing Democratic reluctance to engage in a politically charged debate over guns. The
National Rifle Association said Thursday that it would ``continue to vigorously oppose any
international effort to restrict the constitutional rights of law-abiding American gun owners.''

Since taking office, Mr.~Obama and his aides have been working assiduously to carve out a Mexico
policy that talks of ``shared responsibility'' in combating the drug problem. On Wednesday, the
Obama administration announced stiff financial sanctions against members of three more Mexican drug
cartels, designating them ``kingpins'' under a law that allows the American government to seize
their assets.

In words that resounded on both sides of the border, Secretary of State Hillary Rodham Clinton said
in Mexico City last month that America's ``insatiable demand'' for illegal drugs fueled the trade
and that America's ``inability'' to stop weapons from being smuggled south fed the violence. It is a
marked shift in tone from previous administrations, and Mr.~Obama used his visit here to reiterate
the sentiment.

``I will not pretend that this is Mexico's responsibility alone,'' he said. ``The demand for these
drugs inside the United States is keeping these cartels in business.''

As if to underscore his point, the Mexican authorities reported Thursday that at least 15 drug
cartel gunmen had been killed in a gun battle with army troops. One soldier was also reported killed
in the clash, which took place in San Miguel Totolapan, about 125 miles southwest of Mexico City.

``For the last 30 years the United States has come down with the big sticks of eradication and
helicopters, and the elephant in the room of our own consumption, and the tough proliferation of
arms, were just never addressed,'' said Julia E.~Sweig, director of the Latin America program at the
Council on Foreign Relations. ``I think just beginning to talk about those things is going to buy
him a lot of space down there.''

The drug violence is so intense here that in December, a Pentagon report warned that Mexico could be
on the verge of becoming a failed state. Mr.~Calderón dismissed that assertion in an interview on
Wednesday with the ABC News program ``Nightline.''

Mr.~Obama comes here fresh from a much-publicized swing through Europe that put him squarely on the
world stage. The Latin America trip, which will include a visit to Trinidad and Tobago, may not be
as high profile; for one thing, the first lady, Michelle Obama, who added a touch of pizzazz to the
Europe trip, stayed in Washington. But the president has made repairing relations with world leaders
a signature of his foreign policy, and the visit is designed to give him a chance to do that in a
region with which he is less familiar.

Mr.~Obama was spending the night in Mexico before leaving for Port of Spain, Trinidad, to attend a
gathering of leaders of Western Hemisphere nations.

Cuba is likely to be high on the agenda there. The White House announced this week that Mr.~Obama
was lifting longstanding restrictions on travel and remittances to the island, but some Latin
American leaders want him to do more, and would like to see the United States lift its 47-year-old
trade embargo with Cuba.

Mr.~Obama said he had made a good faith effort to improve relations and was now looking for Cuba to
demonstrate its willingness to make democratic reforms.

\pagebreak
\printindex
\end{document}
